\documentclass[a4paper,11pt]{amsart}



\usepackage{amsmath,amssymb,amsfonts}
\usepackage{amsthm}
\usepackage{xcolor}
\usepackage{geometry}
\usepackage{microtype}
\usepackage{aliascnt}
\usepackage[hidelinks]{hyperref}
\usepackage[nameinlink,capitalise]{cleveref}

\geometry{a4paper,margin=2.7cm}
\allowdisplaybreaks
\numberwithin{equation}{section}

\newtheorem{theorem}{Theorem}[section]
\newaliascnt{proposition}{theorem}
\newtheorem{proposition}[proposition]{Proposition}
\aliascntresetthe{proposition}
\newaliascnt{lemma}{theorem}
\newtheorem{lemma}[lemma]{Lemma}
\aliascntresetthe{lemma}
\newaliascnt{corollary}{theorem}
\newtheorem{corollary}[corollary]{Corollary}
\aliascntresetthe{corollary}

\crefname{theorem}{Theorem}{Theorems}
\crefname{proposition}{Proposition}{Propositions}
\crefname{lemma}{Lemma}{Lemmas}
\crefname{corollary}{Corollary}{Corollaries}

\theoremstyle{definition}
\newtheorem{definition}[theorem]{Definition}
\newaliascnt{example}{theorem} 
\newtheorem{example}[example]{Example} \aliascntresetthe{example}
\theoremstyle{remark}
\newtheorem{remark}[theorem]{Remark}

\crefname{example}{Example}{Examples} 
\Crefname{example}{Example}{Examples}

\usepackage{booktabs,array,longtable}
\setcounter{MaxMatrixCols}{30}
\renewcommand{\arraystretch}{1.12}

\newcommand{\R}{\mathbb R}
\newcommand{\Prob}{\mathbb P}
\newcommand{\Om}{\Omega}
\newcommand{\inter}{\bigcap_{i=1}^n}
\newcommand{\mat}[1]{\begin{pmatrix}#1\end{pmatrix}}



\hypersetup{hypertexnames=false,
 pdftitle={Multivariate copulas with given values at three arbitrary points}}
\title[Multivariate copulas with given values at three points]
{Multivariate copulas with\\given values at three arbitrary points}
% Insert the authors' names, affiliations, and acknowledgments before submission.
\date{}
\keywords{copula, quasi-copula, finite interpolation, checkerboard copula,
Fr\'echet--Hoeffding bounds, double description, Farkas' lemma}

\begin{document}
\begin{abstract}
We characterize the values that a multivariate copula can attain at three
arbitrary points of the unit cube, in every dimension $n\geq2$.
The one-point Fr\'echet--Hoeffding bounds and the directed pairwise
difference bounds must be supplemented by three upper inequalities and
one lower inequality involving all three prescribed values. These four
additional inequalities are redundant in dimension two.
Necessity follows from elementary probability inequalities. Sufficiency
combines exact finite facet calculations with a coordinate-merging lemma
that preserves the criterion in arbitrary dimension. We give an expanded
proof of this lemma, including its simultaneous feasibility formulation,
the application of Farkas' alternative, and the rational projection
certificate. The sufficiency proof is computer-assisted, with reproducible
verification using exact arithmetic. Complete configuration tables are
provided for dimensions two through six, and explicit examples describe
the resulting checkerboard masses. The criterion allows ties,
boundary coordinates, and repeated points; an interpolating checkerboard
copula can be chosen with at most $3n+4$ positive cell masses.
\end{abstract}
\maketitle

\section{Introduction}\label{sec:intro}
A finite interpolation problem asks whether prescribed values of a
multivariate distribution function are compatible with uniform margins.
For two interpolation points, the answer is the same for copulas and
quasi-copulas. This equivalence was established in dimension three by
De Baets, De Meyer, Fern\'andez-S\'anchez, and \'Ubeda-Flores
\cite{DeBaets2013}, who combined Farkas' lemma with the simplex method
to establish feasibility of the associated linear program.
Klement, Kokol Bukov\v{s}ek, Omladi\v{c}, Saminger-Platz, and Stopar
\cite{Klement2023} subsequently established the equivalence in arbitrary
dimension by an analytic argument built on induction on
the dimension. The distinction between the two classes becomes
essential at three points in dimension three, as demonstrated by the
counterexamples in \cite{DeBaets2013}.

Our aim is a complete criterion for copula interpolation at three arbitrary
points in every dimension. Its statement is symmetric in the interpolation
points and does not require a case distinction. In addition to the
quasi-copula conditions, four scalar inequalities suffice, although some
are redundant in a given configuration. A finite mass formulation gives
checkerboard interpolants on the grid generated by the prescribed points.
The examples describe these cell masses and the obstructions that separate
copula interpolation from quasi-copula interpolation.

The necessity proofs are elementary. The sufficiency proof has two
computational ingredients. Exact facet calculations establish the finite
base cases. A dimension-independent projection certificate then proves
that two equally ordered coordinate columns can be merged while
preserving all the interpolation inequalities. There are six possible
orders of three point labels, so such a pair exists in every dimension
greater than six. Conditional gluing restores the two coordinates after
interpolation in the lower dimension. The merging statement, rather than
the coincidence of facet families in a finite list of dimensions, is what
justifies the induction.

The coordinate-merging lemma, \cref{lem:merge-selected}, provides the
dimension-independent reduction underlying the sufficiency proof.
Its proof applies Farkas' lemma to a linear system coupling the reduced
interpolation conditions with those for a bivariate interpolant.
Feasibility follows from a finite family of projection inequalities,
each certified by a nonnegative rational combination of the original
constraints. We provide a complete facet classification in dimension
three and explicit checkerboard constructions illustrating the criterion.
Appendix~\ref{app:all-tables} lists all coordinate-order configurations,
facet counts, and additional facet inequalities in dimensions two through
six. The supplementary programs reproduce the facet calculations and
verify the projection certificates using exact integer and rational
arithmetic.

\section{Finite interpolation and checkerboard extensions}\label{sec:prelim}
An \emph{$n$-copula} is the joint distribution function of $n$ random
variables, each uniform on $[0,1]$. An \emph{$n$-quasi-copula} is a grounded,
coordinatewise nondecreasing map $Q\colon[0,1]^n\to[0,1]$ with uniform margins
and satisfying
\[
 |Q(x)-Q(y)|\leq\sum_{i=1}^n|x_i-y_i|.
\]
Groundedness means that $Q(x)=0$ if some $x_i=0$. Uniform margins mean
\[Q(1,\ldots,1,t,1,\ldots,1)=t.\]
Every copula is a quasi-copula: changing a single threshold from $s$ to
$t\geq s$ changes the corresponding joint probability by at most $t-s$.
We use the componentwise order, so $x\leq y$ means $x_i\leq y_i$ for every
$i$. Componentwise minimum and maximum are denoted by $x\wedge y$ and
$x\vee y$. Put
\begin{align*}
 W_n(x)&=\max\left\{0,\sum_{i=1}^n x_i-n+1\right\},
 &M_n(x)&=\min_i x_i,\\
 D(x,y)&=\sum_{i=1}^n(x_i-y_i)_+,
 &a_+&=\max(a,0).
\end{align*}

The following lemma characterizes finite data admitting a quasi-copula
interpolant in terms of the Fr\'echet--Hoeffding bounds and directed
pairwise inequalities; see \cite[Proposition~3.1]{OVZ}.

\begin{lemma}\label{lem:quasi}
Let $x^1,\ldots,x^k\in[0,1]^n$ and $q_1,\ldots,q_k\in\R$.
There exists a quasi-copula $Q$ with $Q(x^p)=q_p$ for every $p$ if and only if
\begin{align}
 W_n(x^p)&\leq q_p\leq M_n(x^p) &&(1\leq p\leq k),\label{eq:FH}\\
 q_a-q_b&\leq D(x^a,x^b) &&(a\ne b).\label{eq:pair}
\end{align}
\end{lemma}



For later use, an explicit extension under these conditions is
\begin{equation}
 Q(u)=\max\left\{W_n(u),\ \max_{1\leq p\leq k}
             \bigl(q_p-D(x^p,u)\bigr)\right\}.
 \label{eq:quasi-extension}
\end{equation}
Each expression in the maximum is nondecreasing and $1$-Lipschitz in the
$\ell^1$ norm. The one-point upper bounds give $Q(u)\leq M_n(u)$:
for every $i$, $q_p-(x_i^p-u_i)_+\leq u_i$.
Together with $Q\geq W_n$, this yields groundedness and uniform margins.
At $u=x^p$, the directed inequalities bound every term by $q_p$, and the
$p$th term attains that value. 

\subsection{The finite mass problem}
For $k$ prescribed points, choose in each coordinate $i$ a permutation
$\pi_i$ of $\{1,\ldots,k\}$ such that
\[
 x_i^{\pi_i(1)}\leq x_i^{\pi_i(2)}\leq\cdots\leq x_i^{\pi_i(k)}.
\]
Set $t_{i0}=0$, $t_{i,k+1}=1$, and
\[
 t_{ij}=x_i^{\pi_i(j)},\qquad r_{i,\pi_i(j)}=j
 \qquad (j=1,\ldots,k).
\]
Thus $r_{ip}$ records the position assigned to the label $p$ in this
ordering, and $x_i^p=t_{i,r_{ip}}$. These thresholds determine
$k+1$ bands in each coordinate, with widths
\[
 \delta_{ij}=t_{ij}-t_{i,j-1},
 \qquad j=1,\ldots,k+1.
\]
Labels corresponding to equal coordinate values may be ordered
arbitrarily; the bands between successive equal values have length zero.

Let
\[
 J_{ij}=[t_{i,j-1},t_{ij}),\qquad
 B_a=\prod_{i=1}^n J_{i,a_i},\quad a\in\{1,\ldots,k+1\}^n.
\]

\begin{lemma}\label{lem:grid}
The prescribed data admit a copula interpolant if and only if there
are numbers $m_a$, indexed by
$a=(a_1,\ldots,a_n)\in\{1,\ldots,k+1\}^n$, satisfying
\begin{align}
 m_a&\geq0
 &&\bigl(a\in\{1,\ldots,k+1\}^n\bigr),\notag\\
 \sum_{a:a_i=j}m_a&=\delta_{ij}
 &&(1\leq i\leq n,\ 1\leq j\leq k+1),\label{eq:grid-margins}\\
 \sum_{\substack{a:a_i\leq r_{ip}\\
                  \text{for all }i=1,\ldots,n}}m_a&=q_p
 &&(1\leq p\leq k).\label{eq:grid-data}
\end{align}
Both sums range over $a\in\{1,\ldots,k+1\}^n$ subject to the indicated
restrictions. Any such family of masses defines a checkerboard copula
with the required values.
\end{lemma}
\begin{proof}
Given a copula, let $m_a$ be its probability mass in $B_a$. Uniform
margins imply \eqref{eq:grid-margins}; summing the cells below $x^p$
gives \eqref{eq:grid-data}. Cell boundaries have zero probability
since the one-dimensional margins are continuous.

Conversely, summing \eqref{eq:grid-margins} over $j$ for any fixed
$i$ gives $\sum_a m_a=1$. Write
\[
 |B_a|=\prod_{i=1}^n\delta_{i,a_i}
\]
for the $n$-dimensional Lebesgue measure of $B_a$.
If $\delta_{i,a_i}=0$ for some $i$, nonnegativity and the corresponding
marginal equality imply $m_a=0$. Thus every cell of zero volume has
zero mass.

For an interval $J=[s,t)$ of positive length, define
\[
 \ell_J(z)=\min\{1,\max\{0,(z-s)/(t-s)\}\},
\]
the distribution function of the uniform probability measure on $J$.
Distribute each mass $m_a$ uniformly over its cell whenever $|B_a|>0$.
The resulting distribution function and density are
\begin{equation}
\begin{aligned}
 C(z)&=\sum_{a:\,|B_a|>0}
          m_a\prod_{i=1}^n\ell_{J_{i,a_i}}(z_i),\\
 f(z)&=\frac{m_a}{|B_a|}
       \qquad\bigl(z\in\operatorname{int}(B_a),\ |B_a|>0\bigr).
\end{aligned}
\label{eq:checker-formula}
\end{equation}
Set $f=0$ on the grid boundaries; these have Lebesgue measure zero.
Every denominator in \eqref{eq:checker-formula} is therefore positive.

For each band $J_{ij}$ of positive length, the $i$th marginal density
on its interior is
\[
 \frac{\sum_{a:a_i=j}m_a}{\delta_{ij}}=1.
\]
Hence all margins are uniform, so $C$ is a copula.
Finally, \eqref{eq:grid-data} gives $C(x^p)=q_p$ for every $p$.
\end{proof}


\section{A complete three-point criterion in arbitrary dimension}\label{sec:three}
Let $n\geq2$ and $x^1,x^2,x^3\in[0,1]^n$. For distinct $a,b,c\in\{1,2,3\}$ define
\begin{equation}
 T_{(a,b),c}=M_n(x^a\vee x^b)
       +\sum_{i=1}^n\bigl(\min(x_i^a,x_i^b)-x_i^c\bigr)_+.
 \label{eq:T}
\end{equation}
This expression is symmetric in $a,b$. Also put
\begin{equation}
 \ell_i=\min_{p=1,2,3}x_i^p,\qquad
 h_i=\max_{p=1,2,3}x_i^p,\qquad
 B=\sum_{i=1}^n(\ell_i+h_i)-(2n-2).\label{eq:B}
\end{equation}

\begin{theorem}\label{thm:three}
For $n\geq2$, $x^1,x^2,x^3\in[0,1]^n$ and $q_1,q_2,q_3\in\R$, a copula with
$C(x^p)=q_p$ for $p=1,2,3$ exists if and only if the one-point bounds
\eqref{eq:FH}, the directed pairwise bounds \eqref{eq:pair}, and the following
four additional inequalities hold:
\begin{align}
 q_a+q_b-q_c&\leq T_{(a,b), c}
 &&\bigl(\{a,b,c\}=\{1,2,3\}\bigr),\label{eq:triple-upper}\\
 q_1+q_2+q_3&\geq B.\label{eq:triple-lower}
\end{align}
There is one inequality in \eqref{eq:triple-upper} for each choice of $c$.
Equivalently, the data must admit a quasi-copula interpolant and satisfy
\eqref{eq:triple-upper}--\eqref{eq:triple-lower}.
Whenever they do, there is a checkerboard interpolant on the grid generated
by the three points, with at most $3n+4$ cells of positive mass.
For $n=2$, the four additional inequalities are redundant. 
\end{theorem}

The sufficiency proof combines the finite base calculation in
\cref{lem:finite-bases} with coordinate merging and gluing in
\cref{lem:merge-full,lem:merge-glue}. We first give direct probability proofs of all four
additional inequalities. They explain what is missing from pairwise
compatibility.


The following lemma establishes the upper three-point inequalities by
combining inclusion--exclusion with probability bounds obtained from
the uniform margins.

\begin{lemma}\label{lem:upper}
Let $(\Omega,\mathcal F,\Prob)$ be a probability space, with probability
measure $\Prob$, and let $U=(U_1,\ldots,U_n)$ be a random vector on this
space. Assume that each $U_i$ is uniformly distributed on $[0,1]$.
For $x^1,x^2,x^3\in[0,1]^n$, define
\[
 A_p=\bigcap_{i=1}^n
       \{\omega\in\Omega:U_i(\omega)\leq x_i^p\},
 \qquad p=1,2,3.
\]
Then, for distinct $a,b,c\in\{1,2,3\}$,
\[
 \Prob(A_a)+\Prob(A_b)-\Prob(A_c)\leq T_{ab;c}.
\]
\end{lemma}


\begin{proof}
Write $H=A_a\cup A_b$ and $K=A_a\cap A_b$. Inclusion--exclusion gives
\[
 \Prob(A_a)+\Prob(A_b)-\Prob(A_c)
 =\Prob(H)+\Prob(K)-\Prob(A_c)
 \leq\Prob(H)+\Prob(K\setminus A_c).
\]
For every $i$, $H\subseteq\{U_i\leq\max(x_i^a,x_i^b)\}$, so
$\Prob(H)\leq M_n(x^a\vee x^b)$.
Furthermore,
\[
 K\setminus A_c\subseteq
 \bigcup_{i=1}^n\{x_i^c<U_i\leq\min(x_i^a,x_i^b)\}.
\]
The union bound and the uniform margins bound its probability by the sum
in \eqref{eq:T}. These two estimates prove the result.
\end{proof}

The following lemma establishes the lower three-point inequality by
taking expectations of a pointwise bound for the sum of the indicators
of the three lower-orthant events.

\begin{lemma}\label{lem:lower}
Let $\Prob$, $U$, $x^1,x^2,x^3$, and $A_1,A_2,A_3$ be as in
\cref{lem:upper}. Let $\ell_i,h_i$ be as in \eqref{eq:B}. 
Then
\[
 \Prob(A_1)+\Prob(A_2)+\Prob(A_3)
 \geq \sum_{i=1}^n(\ell_i+h_i)-(2n-2).
\]
\end{lemma}
\begin{proof}
Set $L_i=\boldsymbol{1}_{\{U_i\leq\ell_i\}}$ and
$H_i=\boldsymbol{1}_{\{U_i\leq h_i\}}$. We claim the pointwise inequality
\begin{equation}
 \sum_{p=1}^3\boldsymbol{1}_{A_p}
 \geq\sum_{i=1}^n(L_i+H_i)-(2n-2).\label{eq:indicator-lower}
\end{equation}
If some $H_i=0$, then also $L_i=0$ and the right-hand side is nonpositive.
Otherwise all $H_i$ equal one and that side is
$\sum_iL_i-n+2$. It is nonpositive if at most $n-2$ of the $L_i$ equal one.
If exactly $n-1$ equal one, the remaining coordinate lies below its largest
threshold, attained by at least one of the three points. All the other
coordinates lie below every point's thresholds, so at least one event
$A_p$ occurs. Finally, if all $L_i$ equal one, all three events
occur, whereas the right-hand side is two. This proves
\eqref{eq:indicator-lower}. Take expectations and use the uniform margins.
\end{proof}

The upper bounds control the union of two lower orthants and the part of
their intersection outside the third. The lower bound limits how much of
the cube can avoid all three lower orthants while preserving the margins.
In contrast to these constraints, \eqref{eq:pair} considers only two
lower orthants at a time.




\section{Merging equally ordered coordinates}\label{sec:merging}
\subsection{What the lemma is intended to do}
Think of the three interpolation points as the three rows of a matrix.
Two columns, denoted by $u$ and $v$, order these rows in the same way.
After relabeling the rows we may assume
\[
0\le u^1\le u^2\le u^3\le1,
\qquad 0\le v^1\le v^2\le v^3\le1.
\]
Keep one more column $s$ separately, and write the other columns individually
as $y_1,\ldots,y_r$. Thus the original points are
\[
 x^p=(u^p,v^p,s^p,y_1^p,\ldots,y_r^p),\qquad p=1,2,3,
\]
and the original dimension is $r+3$. Every entry is in $[0,1]$.
The column $s$ and each column $y_j$ may have any ordering of the three labels.
The prescribed copula values are $q_1,q_2,q_3$.

The desired operation is to replace the two columns $u,v$ by a single column
$z=(z^1,z^2,z^3)$, producing the points
\[
 \widetilde x^p=(z^p,s^p,y_1^p,\ldots,y_r^p).
\]
The new entries have to do two jobs at once. They must preserve the relevant
interpolation inequalities for the reduced data, and they must themselves
be realizable as the values of a bivariate copula at $(u^p,v^p)$.

% \medskip
% \noindent\textbf{The statement in words.}
% \emph{Assume the original data satisfy the summed interpolation inequalities,
% together with the individual upper bounds associated with the three selected
% columns $u,v,s$. Then there is one ordered triple $z$ that preserves the
% summed inequalities after merging, preserves the individual upper bounds
% associated with $z,s$, and is simultaneously interpolated at $(u^p,v^p)$ by
% a bivariate copula. The number and the orderings of the other columns are
% arbitrary.}

% Here ``associated with a column'' refers both to the one-point upper bound
% and to the upper three-point bound whose outer maximum uses that column.
% The exact inequalities are displayed next.

\subsection{The hypotheses, written with the original coordinates}
The following six families are the inequalities of
\cref{thm:three}, except that individual upper bounds are retained only
for the selected columns $u,v,s$. Empty sums are zero.\\

\noindent \textit{Nonnegativity.}
\begin{equation}
 q_p\ge0
 \qquad(p=1,2,3).
 \tag{H1}\label{eq:H1}
\end{equation}

\noindent \textit{The lower Fr\'echet bound.}
\begin{equation}
 q_p\ge u^p+v^p+s^p+\sum_{j=1}^r y_j^p-(r+2)
 \qquad(p=1,2,3).
 \tag{H2}\label{eq:H2}
\end{equation}
% Equivalently,
% \[
%  q_p-1+(1-u^p)+(1-v^p)+(1-s^p)
%        +\sum_{j=1}^r(1-y_j^p)\ge0
%  \qquad(p=1,2,3).
% \]

\noindent\textit{One-point upper bounds for the coordinates $u,v,s$.}
\begin{equation}
 q_p\le u^p,\qquad q_p\le v^p,\qquad q_p\le s^p
 \qquad(p=1,2,3).
 \tag{H3}\label{eq:H3}
\end{equation}

\noindent \textit{The directed difference bounds.}
\begin{equation}
\begin{aligned}
 q_a-q_b\le{}&
 (u^a-u^b)_++(v^a-v^b)_++(s^a-s^b)_+\\
 &+\sum_{j=1}^r(y_j^a-y_j^b)_+
 \qquad(a,b\in\{1,2,3\},\ a\ne b).
\end{aligned}
\tag{H4}\label{eq:H4}
\end{equation}


\noindent\textit{Upper three-point bounds indexed by $u,v,s$.}
\begin{equation}
\begin{aligned}
 q_a+q_b-q_c\le{}&
 \min\bigl\{\max(u^a,u^b),\max(v^a,v^b),\max(s^a,s^b)\bigr\}\\
 &+(\min(u^a,u^b)-u^c)_+
  +(\min(v^a,v^b)-v^c)_+
 +(\min(s^a,s^b)-s^c)_+\\
 &+\sum_{j=1}^r(\min(y_j^a,y_j^b)-y_j^c)_+,
 \qquad \{a,b,c\}=\{1,2,3\}.
\end{aligned}
\tag{H5}\label{eq:H5}
\end{equation}


\noindent \textit{The lower bound for the sum of the three values.}
\begin{equation}
\begin{aligned}
 q_1+q_2+q_3\ge{}&u^1+u^3+v^1+v^3+\min_p s^p+\max_p s^p\\
 &+\sum_{j=1}^r\bigl(\min_p y_j^p+\max_p y_j^p\bigr)-(2r+4).
\end{aligned}
\tag{H6}\label{eq:H6}
\end{equation}

\noindent\emph{Relation to the full criterion.}
The full interpolation criterion also includes $q_p\le y_j^p$ and the
upper bound (H5) with outer maximum $\max(y_j^a,y_j^b)$, for every $j$.
Those individual bounds are not hypotheses of \cref{lem:merge-selected}. Its role is to
prove a statement with one other selected column $s$. Section~\ref{subsec:merge-common}
explains why this is enough for the full coordinate-merging lemma.


\subsection{The conclusion, also written without grouping}
\begin{lemma}[Selected-coordinate projection]
\label{lem:merge-selected}
Let $r$ be a nonnegative integer, and let
\[
 u=(u^1,u^2,u^3),\quad
 v=(v^1,v^2,v^3),\quad
 s=(s^1,s^2,s^3),\quad
 y_j=(y_j^1,y_j^2,y_j^3)\;\;(1\le j\le r)
\]
be vectors in $[0,1]^3$. Assume that
\[
 0\le u^1\le u^2\le u^3\le1,
 \qquad
 0\le v^1\le v^2\le v^3\le1.
\]
No ordering is imposed on the entries of $s$ or of any $y_j$.
Let $q_1,q_2,q_3\in\mathbb R$ be prescribed values at the points
\[
 x^p=(u^p,v^p,s^p,y_1^p,\ldots,y_r^p)\in[0,1]^{r+3},
 \qquad p=1,2,3,
\]
and suppose that these coordinates and values satisfy
\eqref{eq:H1}--\eqref{eq:H6}.

Then there exist numbers
\[
 0\le z^1\le z^2\le z^3\le1
\]
such that the reduced points
\[
 \widetilde x^p=(z^p,s^p,y_1^p,\ldots,y_r^p)\in[0,1]^{r+2},
 \qquad p=1,2,3,
\]
with the same prescribed values $q_1,q_2,q_3$, satisfy the following
inequalities:
\begin{align}
 q_p&\ge0
 &&(p=1,2,3),\tag{R1}\label{eq:R1}\\
 q_p&\ge z^p+s^p+\sum_{j=1}^r y_j^p-(r+1)
 &&(p=1,2,3),\tag{R2}\label{eq:R2}\\
 q_p&\le z^p,\qquad q_p\le s^p
 &&(p=1,2,3),\tag{R3}\label{eq:R3}
\end{align}
\begin{equation}
\begin{aligned}
 q_a-q_b\le{}&
 (z^a-z^b)_++(s^a-s^b)_+
 +\sum_{j=1}^r(y_j^a-y_j^b)_+\\
 &\hspace{2em}(a,b\in\{1,2,3\},\ a\ne b),
\end{aligned}
\tag{R4}\label{eq:R4}
\end{equation}
\begin{equation}
\begin{aligned}
 q_a+q_b-q_c\le{}&
 \min\bigl\{\max(z^a,z^b),\max(s^a,s^b)\bigr\}\\
 &+(\min(z^a,z^b)-z^c)_+
  +(\min(s^a,s^b)-s^c)_+\\
 &+\sum_{j=1}^r(\min(y_j^a,y_j^b)-y_j^c)_+
 \qquad(\{a,b,c\}=\{1,2,3\}),
\end{aligned}
\tag{R5}\label{eq:R5}
\end{equation}
and
\begin{equation}
\begin{aligned}
 q_1+q_2+q_3\ge{}&
 z^1+z^3+\min_{1\le p\le3}s^p+\max_{1\le p\le3}s^p\\
 &+\sum_{j=1}^r
 \left(\min_{1\le p\le3}y_j^p+\max_{1\le p\le3}y_j^p\right)
 -(2r+2).
\end{aligned}
\tag{R6}\label{eq:R6}
\end{equation}
Empty sums are understood as zero.

Moreover, for this same triple $(z^1,z^2,z^3)$, there exists a
bivariate copula $B$ satisfying
\[
 B(u^p,v^p)=z^p\qquad(p=1,2,3).
\]
\end{lemma}
\subsubsection{The simultaneous bivariate requirement}
For the three equally ordered points $(u^p,v^p)$, the bivariate interpolation criterion is
\begin{align}
 z^p&\ge0,\qquad z^p\ge u^p+v^p-1,\qquad
 z^p\le u^p,\qquad z^p\le v^p,
 \tag{C1}\label{eq:C1}\\
 z^a-z^b&\le(u^a-u^b)_++(v^a-v^b)_+\qquad(a\ne b).
 \tag{C2}\label{eq:C2}
\end{align}
For equally ordered $u,v$, the second line enforces monotonicity of $z$
and bounds its increments. In particular, for $p=1,2$,
\[
0\le z^{p+1}-z^p\le(u^{p+1}-u^p)+(v^{p+1}-v^p).
\]
Conditions (C1)--(C2) suffice by the bivariate facet calculation in
\cref{lem:finite-bases}. That finite calculation is independent of the
coordinate-merging argument. It gives a bivariate checkerboard copula
$B$ with $B(u^p,v^p)=z^p$ for all three prescribed points.

The exact verifier retains the following additional valid, redundant
bivariate inequalities as well. They follow from \cref{lem:upper,lem:lower}
applied to this $B$. 
\begin{equation}
\begin{aligned}
 z^a+z^b-z^c\le{}&
 \min\bigl\{\max(u^a,u^b),\max(v^a,v^b)\bigr\}\\
 &+(\min(u^a,u^b)-u^c)_+
  +(\min(v^a,v^b)-v^c)_+
 \qquad(\{a,b,c\}=\{1,2,3\}).
\end{aligned}
\tag{C3}\label{eq:C3}
\end{equation}
 and
\begin{equation}
 z^1+z^2+z^3\ge u^1+u^3+v^1+v^3-2.
 \tag{C4}\label{eq:C4}
\end{equation}
Including (C3)--(C4) makes the displayed system agree with the one actually
used in the certificate. They do not impose a stronger bivariate interpolation
requirement.

\subsection{The projection proof, step by step}\label{subsec:merge-proof}
\subsubsection{The only unknowns are the three entries of the merged column}
Fix the original data and arrange the entries in each column in
nondecreasing order. Entries with equal values may be ordered arbitrarily. Impose the order $z^1\le z^2\le z^3$. Every minimum, maximum,
and positive part in (R1)--(R6) and (C1)--(C4) then selects a fixed linear
expression. For example,
\[
 \max(z^2,z^3)=z^3,\qquad
 (\min(z^2,z^3)-z^1)_+=z^2-z^1.
\]
There is therefore one simultaneous system of linear inequalities in
$z^1,z^2,z^3$. Include the inequalities
$z^1\ge0$, $z^2-z^1\ge0$, $z^3-z^2\ge0$, and $1-z^3\ge0$ in this system.
Write each row as
\begin{equation}\label{eq:row}
 h_k+b_{k1}z^1+b_{k2}z^2+b_{k3}z^3\ge0.
\end{equation}
The coefficients $b_{kp}$ are fixed integers. The quantity $h_k$ is an affine
expression in the original data; with those data fixed, it is a known number.


\subsubsection{Exactly what Farkas' lemma says here}
Choose nonnegative weights $\lambda_k$ for which
\begin{equation}\label{eq:cancellation}
 \sum_k\lambda_k b_{k1}=\sum_k\lambda_k b_{k2}
 =\sum_k\lambda_k b_{k3}=0.
\end{equation}
Multiplying \eqref{eq:row} by these weights and adding eliminates all three
unknowns. It gives the necessary condition
\begin{equation}\label{eq:projected}
 \sum_k\lambda_k h_k\ge0.
\end{equation}
By Farkas' alternative \cite[Corollary~7.1e]{Schrijver1986},
exactly one of the following holds:
\begin{enumerate}
\item There is a single vector $(z^1,z^2,z^3)\in\R^3$ satisfying every row.
\item There are weights $\lambda_k\ge0$ satisfying \eqref{eq:cancellation}
      for which $\sum_k\lambda_kh_k<0$.
\end{enumerate}
% The order and interval restrictions on $z$ have already been included among
% the rows; there is no separate unstated sign restriction in this alternative.
% Thus proving \eqref{eq:projected} for \emph{all} cancelling nonnegative
% combinations is sufficient for a common feasible $z$.

% This is the logical bridge that an informal argument can miss. Checking a
% few individual upper and lower bounds on $z$ would establish only some
% necessary conditions. The proof must account for every possible cancelling
% combination.

\subsubsection{Why only combinations of at most four nonzero directions matter}
Consider the cone of multipliers
\[
 \mathcal K=\{\lambda:\lambda_k\ge0,\ \sum_k\lambda_k b_k=0\},
 \qquad b_k=(b_{k1},b_{k2},b_{k3}).
\]
It is a pointed polyhedral cone because it lies in the nonnegative orthant.
It is generated by its extreme rays. It therefore suffices to check
\eqref{eq:projected} for those rays.

Write the $k$th inequality as
\[
 h_k+b_{k1}z^1+b_{k2}z^2+b_{k3}z^3\ge0.
\]
If $b_k=(b_{k1},b_{k2},b_{k3})=0$, this inequality already gives
a condition on the original data, and its unit multiplier is treated
separately.

Consider an extreme ray of the cone of nonnegative multipliers
satisfying the cancellation equations. Suppose its support is
$S=\{k_1,\ldots,k_t\}$ and all the corresponding vectors $b_{k_j}$
are nonzero. Form the matrix
\[
 M_S=
 \begin{pmatrix}
 b_{k_1,1}&b_{k_2,1}&\cdots&b_{k_t,1}\\
 b_{k_1,2}&b_{k_2,2}&\cdots&b_{k_t,2}\\
 b_{k_1,3}&b_{k_2,3}&\cdots&b_{k_t,3}
 \end{pmatrix}
 \in\mathbb R^{3\times t}.
\]
Thus each column of $M_S$ contains the three coefficients of one
inequality in the original system. The cancellation equations become
\[
 M_S
 \begin{pmatrix}
 \lambda_{k_1}\\
 \lambda_{k_2}\\
 \vdots\\
 \lambda_{k_t}
 \end{pmatrix}
 =
 \begin{pmatrix}0\\0\\0\end{pmatrix},
 \qquad
 \lambda_{k_j}>0\quad(j=1,\ldots,t).
\]

The nullspace of $M_S$ must be one-dimensional. Indeed, if it had
dimension at least two, we could choose
$\eta\in\ker M_S$ that is not proportional to
$\lambda_S=(\lambda_{k_1},\ldots,\lambda_{k_t})^\top$.
Since every entry of $\lambda_S$ is positive, both
$\lambda_S+\varepsilon\eta$ and $\lambda_S-\varepsilon\eta$
would remain nonnegative for sufficiently small $\varepsilon>0$.
Both satisfy the cancellation equations, and their average is
$\lambda_S$, although neither is proportional to $\lambda_S$.
Extending these vectors by zero outside $S$ would contradict
extremality of the original ray.

Consequently,
\[
 \operatorname{rank}M_S=t-1\le3,
 \qquad\text{and hence}\qquad t\le4.
\]
These minimal positive dependencies are called \emph{positive circuits}.

Moreover, $M_S$ cannot contain two identical nonzero columns:
increasing one of their multipliers and decreasing the other by the
same sufficiently small amount would preserve nonnegativity and
cancellation, again contradicting extremality.
The same argument, with the changes scaled appropriately, excludes
two positively proportional columns.


\subsubsection{How the circuits are found exactly}

For this calculation, an upper bound involving the minimum of several
expressions is first replaced by the equivalent family of separate
upper bounds. After resolving the coordinate orders, write all the
inequalities as
\[
 h_k+b_{k1}z^1+b_{k2}z^2+b_{k3}z^3\ge0.
\]
Here $h_k$ is an affine expression in the original coordinates and
prescribed values, whereas
\[
 b_k=(b_{k1},b_{k2},b_{k3})^\top\in\mathbb Z^3
\]
is a fixed coefficient vector.

\paragraph{Step 1: Choose the coefficient vectors and check the rank.}
Let $\mathcal D$ be the set of distinct nonzero vectors occurring among
the $b_k$. For each $t\in\{2,3,4\}$ and each subset
$\{d^1,\ldots,d^t\}\subseteq\mathcal D$, form
\[
 M=
 \begin{pmatrix}
 d^1_1&d^2_1&\cdots&d^t_1\\
 d^1_2&d^2_2&\cdots&d^t_2\\
 d^1_3&d^2_3&\cdots&d^t_3
 \end{pmatrix}.
\]
We seek strictly positive weights satisfying
\[
 M\lambda=0,
 \qquad
 \lambda=(\lambda_1,\ldots,\lambda_t)^\top,
 \qquad
 \lambda_j>0\quad(j=1,\ldots,t).
\]
For these vectors to form a positive circuit, the preceding argument
shows that
\[
 \operatorname{rank}M=t-1.
\]
We therefore discard the subset if this rank condition fails.

\paragraph{Step 2: Compute a null vector using determinants.}
Choose $t-1$ linearly independent rows of $M$, and denote the resulting
$(t-1)\times t$ matrix by $N$.
For $j=1,\ldots,t$, let $N_{\widehat j}$ be the square matrix obtained
by deleting the $j$th column of $N$. Define
\[
 v_j=(-1)^{j+1}\det N_{\widehat j},
 \qquad j=1,\ldots,t.
\]
The cofactor identity gives
\[
 Nv=0,
 \qquad v=(v_1,\ldots,v_t)^\top.
\]
At least one of these determinants is nonzero because
$\operatorname{rank}N=t-1$. Moreover, the chosen rows span the row
space of $M$, so $Mv=0$ as well. Explicitly, the cancellation
identities are
\[
 \sum_{j=1}^t d^j_\ell v_j=0
 \qquad(\ell=1,2,3).
\]

For example, when $t=3$, write
\[
 N=
 \begin{pmatrix}
 \alpha_1&\alpha_2&\alpha_3\\
 \beta_1&\beta_2&\beta_3
 \end{pmatrix}.
\]
The formula becomes
\[
 v=
 \begin{pmatrix}
 \alpha_2\beta_3-\alpha_3\beta_2\\
 \alpha_3\beta_1-\alpha_1\beta_3\\
 \alpha_1\beta_2-\alpha_2\beta_1
 \end{pmatrix}.
\]
When $t=4$, all three rows of $M$ are used, and
\[
 v=
 \begin{pmatrix}
 \det[d^2\ d^3\ d^4]\\
 -\det[d^1\ d^3\ d^4]\\
 \det[d^1\ d^2\ d^4]\\
 -\det[d^1\ d^2\ d^3]
 \end{pmatrix}.
\]

\paragraph{Step 3: Check positivity and normalize the weights.}
Since $\ker M$ is one-dimensional, every cancellation vector supported
on this subset is a scalar multiple of $v$. Consequently, strictly
positive weights exist precisely when
\[
 v_j>0\quad\text{for all }j,
 \qquad\text{or}\qquad
 v_j<0\quad\text{for all }j.
\]
If an entry is zero, no vector in $\ker M$ is strictly positive on
the whole subset. If the entries have mixed signs, no scalar multiple
of $v$ is strictly positive. In either case, the subset is discarded.

Otherwise, choose $\sigma\in\{-1,1\}$ so that $\sigma v_j>0$ for
every $j$, and set
\[
 g=\gcd(|v_1|,\ldots,|v_t|),
 \qquad
 \lambda_j=\frac{\sigma v_j}{g}
 \quad(j=1,\ldots,t).
\]
Then
\[
 \lambda_j\in\mathbb Z_{>0},\qquad
 \gcd(\lambda_1,\ldots,\lambda_t)=1,\qquad
 \sum_{j=1}^t\lambda_jd^j=0.
\]
Thus $\lambda$ is the primitive positive integer cancellation vector
for this circuit.

For instance, consider
\[
 d^1=
 \begin{pmatrix}1\\0\\0\end{pmatrix},
 \qquad
 d^2=
 \begin{pmatrix}0\\1\\0\end{pmatrix},
 \qquad
 d^3=
 \begin{pmatrix}-1\\-1\\0\end{pmatrix}.
\]
Here
\[
 M=
 \begin{pmatrix}
 1&0&-1\\
 0&1&-1\\
 0&0&0
 \end{pmatrix},
 \qquad
 N=
 \begin{pmatrix}
 1&0&-1\\
 0&1&-1
 \end{pmatrix}.
\]
The three signed minors give
\[
 v_1=1,\qquad v_2=1,\qquad v_3=1.
\]
Hence the primitive cancellation vector is $(1,1,1)^\top$, expressing
the identity $d^1+d^2+d^3=0$.

\paragraph{Step 4: Use every input inequality with the required vectors.}
The circuit calculation depends only on the coefficient vectors.
Different inequalities may have the same coefficient vector but
different affine expressions $h_k$. To account for all of them, define
\[
 \mathcal I(d)=\{k:b_k=d\}.
\]
For a circuit with vectors $d^1,\ldots,d^t$ and weights
$\lambda_1,\ldots,\lambda_t$, consider every tuple
\[
 (k_1,\ldots,k_t)
 \in
 \mathcal I(d^1)\times\cdots\times\mathcal I(d^t).
\]
Multiplying the selected inequalities by their circuit weights and
adding them gives
\[
 \sum_{j=1}^t\lambda_j
 \left(h_{k_j}+\sum_{\ell=1}^3d^j_\ell z^\ell\right)
 =
 \sum_{j=1}^t\lambda_jh_{k_j}\ge0,
\]
because the coefficients of all three variables $z^\ell$ cancel.
Thus this choice produces the projected inequality
\[
 \sum_{j=1}^t\lambda_jh_{k_j}\ge0.
\]

For example, suppose the input contains
\[
 h_1+z^1\ge0,\qquad
 h'_1+z^1\ge0,\qquad
 h_2+z^2\ge0,\qquad
 h_3-z^1-z^2\ge0.
\]
The circuit $(1,1,1)^\top$ above produces both inequalities
\[
 h_1+h_2+h_3\ge0,
 \qquad
 h'_1+h_2+h_3\ge0.
\]
Keeping only one representative of the coefficient vector
$(1,0,0)^\top$ would omit one of these conditions.

\paragraph{Step 5: Include zero vectors and remove repetitions.}
For every input inequality with $b_k=0$, include the condition
\[
 h_k\ge0.
\]
Together with the circuit inequalities, these account for all extreme
rays of the multiplier cone.

Finally, expand each projected inequality as an affine expression in
the original data. Clear denominators by a positive common
denominator and divide its integer coefficients, including the
constant term, by their greatest common divisor.
Only positive rescaling is used, so the direction of the inequality
is preserved. Remove identically zero inequalities and repeated
normalized coefficient vectors.

By Farkas' alternative, the resulting finite list characterizes
feasibility of the system in $z$. To complete the projection proof,
one must then verify that every inequality in this list follows from
\eqref{eq:H1}--\eqref{eq:H6} and the coordinate-order assumptions.




\subsection{Two complete examples of elimination}\label{subsec:merge-examples}
\subsubsection{A two-row cancellation recovers the original lower bound}
Take the reduced lower bound (R2) and the local lower bound from (C1):
\begin{align*}
 q_p-z^p-s^p-\sum_{j=1}^r y_j^p+r+1&\ge0,\\
 z^p-u^p-v^p+1&\ge0.
\end{align*}
Their $z$-coefficient vectors are $-e_p$ and $e_p$. Adding with weights one
cancels $z^p$ and gives
\[
 q_p-u^p-v^p-s^p-\sum_{j=1}^r y_j^p+r+2\ge0,
\]
which is exactly hypothesis (H2). Thus this possible obstruction is excluded
by an original assumption.

\subsubsection{A four-row cancellation uses all three unknowns}
Use (R5) with $(a,b,c)=(2,3,1)$ and outer column $w=z$, together with three
local one-point bounds. The four inequalities are
\begin{align*}
0\le{}&q_1-q_2-q_3+z^3+z^2-z^1
 +(\min(s^2,s^3)-s^1)_+
 +\sum_{j=1}^r(\min(y_j^2,y_j^3)-y_j^1)_+,\\
0\le{}&u^2-z^2,\\
0\le{}&v^3-z^3,\\
0\le{}&z^1-u^1-v^1+1.
\end{align*}
The corresponding vectors are
\[
 (-1,1,1),\qquad(0,-1,0),\qquad(0,0,-1),\qquad(1,0,0).
\]
They sum to zero with all weights equal to one; no proper subcollection has
a positive cancellation. Adding the four inequalities gives $\rho\ge0$, where
\begin{align*}
\rho={}&q_1-q_2-q_3+u^2+v^3-u^1-v^1+1
 +(\min(s^2,s^3)-s^1)_++\sum_{j=1}^r(\min(y_j^2,y_j^3)-y_j^1)_+.
\end{align*}
Why must this expression be nonnegative under the hypotheses? There is an
explicit identity:
\begin{align*}
\rho={}&\Bigl[q_1-q_2-q_3+v^3+(u^2-u^1)+(v^2-v^1)\\
 &\qquad +(\min(s^2,s^3)-s^1)_+
 +\sum_{j=1}^r(\min(y_j^2,y_j^3)-y_j^1)_+\Bigr]\\
 &+(v^3-v^2)+(1-v^3).
\end{align*}
The bracket is the nonnegative slack of (H5) with outer column $v$.
The last two terms are nonnegative ordered-grid widths. This is a complete
Farkas certificate: the projected inequality is a sum of original
inequalities with nonnegative coefficients. The supplied supplementary
check verifies this identity against the actual input and hypothesis rows
for all six orders of $s$.

\subsection{How the finite certificate covers arbitrarily many columns}
The expanded inequalities above contain every $y_j$ separately. To perform
one finite calculation for every $r$, the implementation records the sums
that occur in them using a fixed number of variables. This is an encoding
step in the verification, rather than a replacement for their interpretation.

\subsubsection{The exact identities behind grouping}
For example, if $I$ is a set of columns sharing one order of the three labels,
then
\[
 \sum_{j\in I}(y_j^a-y_j^b)_+
 =\left(\sum_{j\in I}y_j^a-\sum_{j\in I}y_j^b\right)_+.
\]
All differences have the same weak sign. Also the same label attains the
minimum, or maximum, in every column of the group, so
\begin{align*}
 \sum_{j\in I}\min_p y_j^p&=\min_p\sum_{j\in I}y_j^p,\\
 \sum_{j\in I}\max_p y_j^p&=\max_p\sum_{j\in I}y_j^p,\\
 \sum_{j\in I}(\min(y_j^a,y_j^b)-y_j^c)_+
 &=\left(\min\left\{\sum_{j\in I}y_j^a,\sum_{j\in I}y_j^b\right\}
                     -\sum_{j\in I}y_j^c\right)_+.
\end{align*}
The lower-bound deficit sums satisfy the immediate identity
$\sum_{j\in I}(1-y_j^p)=|I|-\sum_{j\in I}y_j^p$.
These are exact identities, not estimates.

\subsubsection{Why arbitrary nonnegative band totals are allowed}
A column ordered as $y_j^{\tau(1)}\le y_j^{\tau(2)}\le y_j^{\tau(3)}$
has four nonnegative widths
\[
 y_j^{\tau(1)},\quad y_j^{\tau(2)}-y_j^{\tau(1)},\quad
 y_j^{\tau(3)}-y_j^{\tau(2)},\quad1-y_j^{\tau(3)}.
\]
Sum each width over all columns having that order. There are six possible
orders, giving 24 nonnegative band totals. Their values encode every summed
term in (H1)--(H6) and (R1)--(R6), regardless of $r$.

The exact certificate also proves a stronger version of
\cref{lem:merge-selected}: these 24 totals may be \emph{arbitrary} nonnegative real numbers.
An equivalent ungrouped way to state this stronger assertion is to allow each
remaining column $y_j$ a nonnegative weight $w_j$. In the displayed
hypotheses and conclusions replace
\[
 \sum_{j=1}^r f(y_j)\quad\text{by}\quad\sum_{j=1}^r w_j f(y_j),
 \qquad r\quad\text{by}\quad\sum_{j=1}^r w_j,
\]
for each summed expression, with the replacement of $r$ made in the
lower-bound constants only. The same merging conclusion
holds. The usual coordinate statement takes every $w_j=1$.

To see that this weighted wording includes arbitrary totals, take the four
totals of one order and add them to obtain a number $w$. If $w>0$, divide
those totals by $w$; they are the four band widths of one column in $[0,1]^3$,
which is given weight $w$. If $w=0$, omit that group. Conversely, summing
weighted column widths gives arbitrary nonnegative band totals of exactly
the type used by the certificate. The weighted extension is an algebraic
statement; noninteger weights are not interpreted as a fractional number of
copula coordinates.

\subsubsection{What is checked, and why it completes the proof}
For each order of $s$, the calculation retains three free band widths for each
of $u,v,s$, the 24 band totals, and $q_1,q_2,q_3$: $9+24+3=36$ variables.
An extra entry equal to one records constants. Appending the three unknown
entries of $z$ makes each conclusion row an integer vector of length 40.
The calculation reconstructs the following counts:
\begin{center}
\begin{tabular}{lr}
\toprule
Quantity, for one order of $s$ & Count\\
\midrule
Original hypothesis and grid rows & 67\\
Combined reduced and local rows, after deduplication & 87\\
Distinct nonzero coefficient vectors in $z$ & 18\\
Positive circuits among these vectors & 160\\
Distinct projected inequalities, including zero-$z$ rows & 1,893\\
\bottomrule
\end{tabular}
\end{center}
For every projected affine expression $\rho$, the certificate supplies an
identity
\begin{equation}\label{eq:identity}
 \rho=\sum_{\ell=1}^{67}\mu_\ell\varphi_\ell,
 \qquad \mu_\ell\in\mathbb Q,\quad\mu_\ell\ge0,
\end{equation}
where $\varphi_\ell\ge0$ is one of the original hypotheses or ordered-grid
inequalities, represented in those retained variables. Section~\ref{subsec:merge-examples}
gives two instances of this reasoning in ordinary coordinate notation.
The verifier checks equality of all 37 retained entries, including the
constant entry, in every identity \eqref{eq:identity}.

All hypotheses make every $\varphi_\ell$ nonnegative, so every projected
inequality is nonnegative. The circuit enumeration is complete by the
extreme-ray argument in Section~\ref{subsec:merge-proof}; hence all the multipliers
required by Farkas' alternative are covered. That alternative supplies one
common ordered triple $z$ satisfying every reduced and local row. The
bivariate criterion then supplies the copula $B$. This proves \cref{lem:merge-selected},
including its stronger weighted version.

There are six possible orders of $s$. Each produces 1,893 certified
expressions, so the verifier checks $6\cdot1,893=11,358$ exact identities.
It uses integer and rational arithmetic, with no floating-point tolerance
or optimization solver during verification. The finite certificate check is
an essential computational part of this proof; the Farkas argument alone
does not assert that the identities exist. The supplementary verification log records a successful check of all
six cases.

\subsection{One merged triple for all remaining columns}
\label{subsec:merge-common}
\begin{lemma}[Full coordinate merging]\label{lem:merge-full}
Suppose $n\geq3$ and the data satisfy the inequalities of
\cref{thm:three}. If two coordinate columns $u,v$ have the same weak
ordering of the three point labels, they can be replaced by a single
column $z$ with that ordering so that the reduced data satisfy all the
same inequalities in dimension $n-1$. Simultaneously, a bivariate copula
$B$ can be chosen with $B(u^p,v^p)=z^p$ for $p=1,2,3$.
\end{lemma}
\begin{proof}
Relabel the points so that $u,v$ both have order $123$, and write the
remaining columns as $s,y_1,\ldots,y_r$, where $r=n-3$.
The full coordinate-merging lemma requires the reduced criterion with upper
bounds indexed by \emph{every} remaining column, not only a selected $s$.
It is not enough to choose a different feasible $z$ for each such column.
Instead, form that entire system first and apply the same Farkas projection.

Rows whose $z$-coefficient is zero are already implied by the original full
criterion. For instance, the remaining one-point upper bounds are unchanged.
In a directed bound with $a<b$, the contributions from $u,v,z$ are all zero.
In an upper three-point inequality with $c=2$ or $c=3$, the leakage from each
of these ordered columns is also zero, so a bound using a remaining column
as its outer maximum is unchanged.

Among rows with a nonzero $z$-coefficient, the only ones indexed by a choice
of a remaining column $w$ have $(a,b,c)=(2,3,1)$ and the form
\begin{align*}
0\le{}&q_1-q_2-q_3+\max(w^2,w^3)+(z^2-z^1)\\
 &+(\min(s^2,s^3)-s^1)_+
 +\sum_{j=1}^r(\min(y_j^2,y_j^3)-y_j^1)_+.
\end{align*}
Here $w$ ranges over $s,y_1,\ldots,y_r$; the sum of positive parts always
includes all remaining columns, including the one used for the outer
maximum. Every such row has the same coefficient vector $(-1,1,0)$.
A positive circuit can contain at most one of them.

Take any circuit of the full system. If it uses one of these indexed rows,
choose that row's column as the selected column $s$ for applying \cref{lem:merge-selected};
otherwise choose any remaining column. All the rows in this circuit now
belong to the system covered by that lemma. Its projected inequality
therefore follows from the original assumptions. Do this for every circuit.
Together with the zero-coefficient rows, all projected inequalities of the
\emph{full} system hold. Farkas' alternative then gives a single $z$ that
satisfies the full system simultaneously.

The selected column may change from one circuit check to another.
No value of $z$ is chosen during those checks. The common $z$ is obtained
only at the end, after all possible obstructions have been excluded. This
is the precise connection between \cref{lem:merge-selected} and the full merging lemma.

\end{proof}

\begin{lemma}[Gluing the merged construction]\label{lem:merge-glue}
If the reduced data of \cref{lem:merge-full} admit a copula interpolant,
then that copula and the bivariate interpolant $B$ can be combined to
obtain an interpolant for the original data.
\end{lemma}
\begin{proof}
Relabel the points so $u^1\le u^2\le u^3$ and
$v^1\le v^2\le v^3$, hence $z^1\le z^2\le z^3$.
For the bivariate interpolant $(S,T)$, the events
$G_p=\{S\le u^p,T\le v^p\}$ are nested and have probabilities $z^p$.
They partition its probability space into four layers of masses
$z^1,z^2-z^1,z^3-z^2,1-z^3$.

Let $(V,U_3,\ldots,U_n)$ realize the lower-dimensional interpolant.
The uniform variable $V$ has the corresponding four interval layers
cut at $z^1,z^2,z^3$. Conditional on its layer, sample $(S,T)$ from the
bivariate interpolant conditioned on the matching $G$-layer,
independently of the other coordinates given that layer. Layers of
probability zero are omitted.
The pair $(S,T)$ retains its bivariate distribution, so both new margins
are uniform. Moreover, $G_p$ now agrees almost surely with
$\{V\le z^p\}$. The prescribed intersections with the remaining
coordinates therefore keep their probabilities $q_p$.
\end{proof}

\subsection{A numerical illustration of the merging operation}
Consider the four-dimensional data
\begin{center}
\begin{tabular}{c|rrrr|r|r}
$p$ & $u^p$ & $v^p$ & $s^p$ & $y_1^p$ & $q_p$ & $z^p$\\
\hline
1 & $.2$ & $.3$ & $.4$ & $.9$ & $.0216$ & $.06$\\
2 & $.5$ & $.6$ & $.7$ & $.6$ & $.126$ & $.30$\\
3 & $.8$ & $.9$ & $.85$ & $.4$ & $.2448$ & $.72$
\end{tabular}
\end{center}
Here $u,v$ have order $123$, whereas $y_1$ has order $321$.
The values were chosen from the product copula:
\[
q_p=u^pv^ps^py_1^p.
\]
Choose $z^p=u^pv^p$. The bivariate product copula gives
$B(u^p,v^p)=z^p$, and the reduced product copula gives
\[
 \widetilde C(z^p,s^p,y_1^p)=z^ps^py_1^p=q_p.
\]
Thus one explicit ordered triple satisfies both requirements.
The entries $z^p$ are probabilities for the two-coordinate events being
merged. They are not obtained by summing the original thresholds. The
product choice works for this example; the general lemma finds a feasible
$z$ without assuming independence.

\section{Finite facet calculations and sufficiency}\label{sec:configurations}
We first give the complete trivariate calculation, with $n=k=3$. Retain the four bands in every coordinate:
\[
 0=t_{i0}\leq t_{i1}\leq t_{i2}\leq t_{i3}\leq t_{i4}=1,
 \qquad \delta_{ij}=t_{ij}-t_{i,j-1}.
\]
An order string such as $132$ means
$x_i^1\leq x_i^3\leq x_i^2$ in the specified coordinate, with these ranks
chosen in the event of a tie. Thus a configuration is a triple of
permutations of $123$. Coordinate permutations and a common relabeling of
the three points preserve interpolation. The $6^3=216$ order triples form
the ten orbits in \cref{tab:counts}.

\subsection{The trivariate data polytopes}
For a fixed rank matrix $R=(r_{ip})$, let $v_a\in\R^{12}$, for
$a\in\{1,2,3,4\}^3$, have entries
\begin{equation}
 v_a=\left(
   (\boldsymbol{1}_{\{a_i=j\}})_{\substack{1\leq i\leq3\\1\leq j\leq3}},
   (\boldsymbol{1}_{\{a_i\leq r_{ip}\ \forall i\}})_{1\leq p\leq3}
       \right).\label{eq:vertices}
\end{equation}
Order the first nine entries by coordinate $i$ and then band $j$, and put
\begin{equation}
 \xi=(\delta_{11},\delta_{12},\delta_{13},
       \delta_{21},\delta_{22},\delta_{23},
       \delta_{31},\delta_{32},\delta_{33},q_1,q_2,q_3).
 \label{eq:z}
\end{equation}
The fourth marginal width is $1-\delta_{i1}-\delta_{i2}-\delta_{i3}$.
By \cref{lem:grid}, the feasible data are exactly
\begin{equation}
 P_R=\operatorname{conv}\{v_a:a\in\{1,2,3,4\}^3\};
 \qquad \xi\in P_R\ \Longleftrightarrow\ \text{copula interpolation}.
 \label{eq:polytope}
\end{equation}
In particular the mass equations use only $13$ rows: total mass, nine
band margins, and three interpolation values. Eliminating the fourth
band from each margin removes the marginal equalities that otherwise
produce lineality in the dual description.
All dual coefficients below are unrestricted real numbers; only the
primal cell masses are constrained to be nonnegative.

\begin{samepage}
\begin{lemma}[Complete facet calculation; computer-assisted]\label{lem:facets}
For each of the ten configurations in \cref{tab:counts}, $P_R$ has affine
dimension $12$. Every facet inequality is one of the following:
\begin{enumerate}
\item a grid inequality $\delta_{ij}\geq0$ $(1\leq j\leq4)$;
\item a linear constituent of \eqref{eq:FH} or a directed inequality
      \eqref{eq:pair};
\item one of the additional inequalities listed in \cref{tab:extras}.
\end{enumerate}
Every listed additional inequality is a facet, after expanding each
minimum of two expressions into two inequalities.
Every inequality in (3) is a linear constituent of
\eqref{eq:triple-upper} or is \eqref{eq:triple-lower}.
The facet numbers are those in \cref{tab:counts}.
\end{lemma}
\end{samepage}

In the tables, \emph{baseline} means the inequalities in (1) and (2),
and \emph{additional} counts facets outside that baseline family.
The facet counts concern the data polytope with variable widths and
values, rather than a fixed-grid slice.
For \cref{tab:extras}, abbreviate
\[
 d_i=t_{i2}-t_{i1},\quad
 U_{ab\mid c}=q_a+q_b-q_c,\quad S=q_1+q_2+q_3,
 \quad B=\sum_i(t_{i1}+t_{i3})-4.
\]

\begin{table}[htbp]
\centering\small
\caption{All order configurations and their exact facet counts.}
\label{tab:counts}
\begin{tabular}{clrrrr}
\toprule
Case & Coordinate orders & Orbit size & Facets & Baseline & Additional\\
\midrule
1 & $(123,123,123)$ & 6  & 24 & 24 & 0\\
2 & $(123,123,132)$ & 18 & 28 & 26 & 2\\
3 & $(123,123,213)$ & 18 & 26 & 26 & 0\\
4 & $(123,123,231)$ & 18 & 29 & 27 & 2\\
5 & $(123,123,312)$ & 18 & 28 & 26 & 2\\
6 & $(123,123,321)$ & 18 & 28 & 26 & 2\\
7 & $(123,132,213)$ & 36 & 28 & 26 & 2\\
8 & $(123,132,231)$ & 36 & 32 & 28 & 4\\
9 & $(123,213,312)$ & 36 & 29 & 27 & 2\\
10& $(123,231,312)$ & 12 & 31 & 27 & 4\\
\bottomrule
\end{tabular}
\end{table}

\begin{table}[htbp]
\centering\small
\caption{Necessary additional inequalities beyond the baseline, by configuration.}
\label{tab:extras}
\begin{tabular}{cp{11.6cm}}
\toprule
Case & Additional inequalities\\
\midrule
1 & None.\\
2 & $U_{23\mid1}\leq\min(t_{13},t_{23})+d_1+d_2+d_3$.\\
3 & None.\\
4 & $U_{13\mid2}\leq\min(t_{13},t_{23})+d_3$.\\
5 & $U_{23\mid1}\leq\min(t_{13},t_{23})+d_1+d_2$.\\
6 & $U_{13\mid2}\leq\min(t_{13},t_{23})$.\\
7 & $U_{23\mid1}\leq\min(t_{13},t_{33})+d_1+d_2$.\\
8 & $U_{13\mid2}\leq t_{22}+d_3$,\quad
    $U_{23\mid1}\leq t_{32}+d_1+d_2$,\newline
    $U_{12\mid3}\leq t_{12}$,\quad $S\geq B$.\\
9 & $U_{23\mid1}\leq\min(t_{13},t_{23})+d_1$.\\
10& $U_{13\mid2}\leq t_{32}+d_2$,\quad
    $U_{23\mid1}\leq t_{22}+d_1$,\newline
    $U_{12\mid3}\leq t_{12}+d_3$,\quad $S\geq B$.\\
\bottomrule
\end{tabular}
\end{table}

\begin{proof}[Exact verification and completeness of the calculation]
We give the finite algorithm used to prove the lemma; its full implementation
and output are supplied with the manuscript.
First generate all triples of permutations of $123$. For each triple,
apply every common permutation of the point labels, sort the three rows
lexicographically, and retain the least result. This gives exactly the
ten representatives and orbit sizes in \cref{tab:counts}; the orbit sizes
sum to $216$.

For each representative construct the $64$ vectors in \eqref{eq:vertices}
and write $h_a=(v_a,1)\in\mathbb Z^{13}$. Exact row reduction gives rank
$13$, so $P_R$ has affine dimension $12$. Its valid affine inequalities
form the cone
\begin{equation}
 K_R=\{y\in\R^{13}:h_a\cdot y\geq0\text{ for all }a\}.
 \label{eq:dual-cone}
\end{equation}
This cone is pointed because the $h_a$ span $\R^{13}$, and it is
full-dimensional because $(0,\ldots,0,1)$ satisfies all inequalities
strictly. The extreme rays of $K_R$ correspond precisely to the facets
of $P_R$, with a ray $y$ giving the inequality $(\xi,1)\cdot y\geq0$.

Use the lexicographic cell order, select the first $13$ linearly
independent rows, and let $H_0$ be the resulting square matrix. The cone
$H_0y\geq0$ is simplicial and its extreme rays are the columns of
$H_0^{-1}$. Compute these columns rationally and rescale each to a
primitive integer vector, preserving its sign. Add the remaining
$51$ inequalities one at a time. The following exact
update is the double-description step \cite{FukudaProdon1996}.

For a new row $h$, retain the old extreme rays $r$ with $h\cdot r\geq0$.
For each adjacent pair $r_+,r_-$ with $h\cdot r_+>0$ and
$h\cdot r_-<0$, add
\begin{equation}
 r_{\mathrm{new}}=(h\cdot r_+)r_--(h\cdot r_-)r_+,
 \label{eq:dd-update}
\end{equation}
again normalized to a primitive integer vector. Deduplicate the result.
Adjacency is determined from the already processed inequalities: if
$Z(r)$ is their set of zero indices at $r$, two distinct rays are adjacent
exactly when no third current ray has a zero set containing
$Z(r_+)\cap Z(r_-)$. Indeed this intersection defines their smallest
common face, and it is two-dimensional precisely when it has just these
two extreme rays. The implementation also uses the necessary preliminary
check $|Z(r_+)\cap Z(r_-)|\geq11$.

These operations enumerate \emph{all} extreme rays, not just a collection
of valid ones. Every new ray created by a half-space cut is the
intersection of its boundary with a two-dimensional face whose two rays
have opposite signs; rays retained from the old cone remain extreme.
The initial list is complete, and this observation proves completeness
at every step by induction. Each intermediate cone remains pointed and
full-dimensional: it is a subcone of the initial pointed cone and still
contains the strictly feasible constant vector above.

Finally, express all proposed inequalities as primitive integer vectors
in the coordinates $(\xi,1)$. The program checks, exactly, that
\begin{enumerate}
\item every enumerated ray is nonnegative on all $64$ input vectors and
      its zero rows have rank $12$;
\item every proposed baseline or triple inequality is nonnegative on
      all $64$ input vectors;
\item every enumerated ray belongs to the baseline or triple family;
\item the enumerated rays outside the baseline family are exactly those
      in \cref{tab:extras}, with the counts in \cref{tab:counts}.
\end{enumerate}
The enumeration uses at most $44$ rays at any intermediate step.
The supplied certificate records every final primitive facet normal,
while the verification program independently rebuilds the complete ray
list before comparing it with the certificate. All matrix inversions,
rank computations, sign tests, and ray updates use integers or rational
numbers. This establishes the finite assertions of the lemma.
\end{proof}

\subsection{The base cases for arbitrary dimension}


For a rank matrix with $n$ coordinate rows, the same construction uses
\begin{equation}
 v_a=\left((\boldsymbol{1}_{\{a_i=j\}})_{\substack{1\leq i\leq n\\1\leq j\leq3}},
       (\boldsymbol{1}_{\{a_i\leq r_{ip}\ \forall i\}})_{p=1,2,3}\right)
       \in\R^{3n+3},\qquad a\in\{1,2,3,4\}^n.
 \label{eq:general-vertices}
\end{equation}
The first $3n$ entries are the first three band margins, and the last
three are the prescribed values. Put $P_R=\operatorname{conv}\{v_a\}$.
By \cref{lem:grid}, interpolation is equivalent to membership of the data
vector $\xi=(\delta_{11},\delta_{12},\delta_{13},\ldots,
\delta_{n1},\delta_{n2},\delta_{n3},q_1,q_2,q_3)$ in $P_R$.

\begin{lemma}[Exact base certificates; computer-assisted]\label{lem:finite-bases}
In dimension two, every facet of $P_R$ is a grid inequality or a
linear constituent of \eqref{eq:FH}--\eqref{eq:pair}.
In dimension six, every facet is a grid inequality or a linear
constituent of \eqref{eq:FH}, \eqref{eq:pair}, \eqref{eq:triple-upper}, or \eqref{eq:triple-lower}.
The exact enumeration gives
\[
\begin{array}{c|rrrrrr}
 n&\text{Order lists}&\text{Types}&\text{Baseline}&\text{Upper}
   &\text{Lower}&\text{Total}\\ \hline
 2&36&5&83&0&0&83\\
 6&46656&83&3865&874&42&4781.
\end{array}
\]
Here upper and lower counts exclude facets already in the baseline
family of grid, one-point, and directed pairwise inequalities.
\end{lemma}

\begin{proof}[Exact enumeration and its completeness]
An order string records one of the six permutations of the three point
labels. Coordinate permutations reduce order lists to multisets of
strings, and simultaneous point relabeling gives the stated symmetry
types. The verifier enumerates all multisets and their label orbits;
the multinomial orbit sizes sum to $6^n$.

For each type, form the rows $h_a=(v_a,1)$ and the cone
$K_R=\{y:h_a\cdot y\ge0\ \forall a\}$. Exact row reduction gives rank
$3n+4$. Thus $P_R$ has dimension $3n+3$ and $K_R$ is pointed. The vector
with only its last coordinate equal to one is strictly feasible, so
the cone is full-dimensional. Its extreme rays are the facet normals
of $P_R$.

The verifier uses exact double description. Start with the first
$3n+4$ independent rows, whose inverse columns generate the initial
simplicial cone. Add the remaining rows one by one. Retain rays on the
nonnegative side of the new row $h$; for every adjacent crossing pair
$r_+,r_-$, add the primitive normalization of
\[
 (h\cdot r_+)r_--(h\cdot r_-)r_+.
\]
Adjacency is checked by common active constraints: no third ray may
have an active set containing their intersection. This tests whether
the smallest face containing the pair is two-dimensional. Every new
extreme ray arises by cutting such a face, so induction proves
completeness of the enumeration.

For every final ray, validity on all $4^n$ vertices and active-row rank
$3n+3$ are checked exactly. Every proposed normal is also checked on
every vertex, and every enumerated normal is matched to the stated
families. Minima are expanded into their individual linear upper
bounds; coordinate orders resolve every positive part and maximum.
The full primitive normal lists are stored in the base certificate.
The verifier reconstructs those lists, rather than merely checking
the stored normals for validity.
\end{proof}

Dimension-six sufficiency also implies sufficiency in every smaller
dimension: append coordinates equal to one to the three points, keep
the prescribed values, and apply the six-dimensional result. All four
families reduce exactly to the original dimension. Restricting the
resulting copula to the face with the appended coordinates equal to
one gives the desired marginal copula. The first assertion of
Lemma~\ref{lem:finite-bases} additionally establishes the bivariate redundancy.


\subsection{Sufficiency in every dimension}
\begin{proof}[Proof of \cref{thm:three}]
Necessity of the one-point and directed pairwise bounds follows from
\cref{lem:quasi}; necessity of the four additional inequalities follows
from \cref{lem:upper,lem:lower}.

For sufficiency, \cref{lem:finite-bases} and the padding argument give the
result for $2\leq n\leq6$. More explicitly, the assumed inequalities and
the nonnegative grid widths satisfy every facet inequality of the relevant
data polytope, so the data vector belongs to it. A convex representation
in its defining vectors supplies the checkerboard masses of \cref{lem:grid}.
For $n=3$, this is also the direct calculation of \cref{lem:facets}:
after permuting coordinates and point labels, the rank configuration is
one of the ten cases in \cref{tab:counts}, and the four symmetric
inequalities contain every additional facet in \cref{tab:extras}.

Suppose now that $n\geq7$ and sufficiency holds in dimension $n-1$.
Choose an ordering of the three point labels in each coordinate,
breaking ties arbitrarily. There are only six possible order strings,
so two coordinates have the same ordering. By \cref{lem:merge-full},
there is a merged column $z$ for which all the reduced inequalities hold
and a bivariate interpolant restoring the two original columns.
The induction hypothesis supplies the reduced copula. The conditional
construction of \cref{lem:merge-glue} gives an interpolant in dimension $n$.
This completes the induction.

The argument includes ties directly. Choose any rank ordering compatible
with the weak coordinate orders. Zero band widths force incident cell
masses to vanish; the checkerboard formula uses only positive-volume cells.
For repeated interpolation points, the two directed bounds force the
prescribed values to agree. No perturbation or limiting argument is needed.

Finally, the data vector lies in the convex hull of the vectors
\eqref{eq:general-vertices} in $\R^{3n+3}$. Carath\'eodory's theorem gives
a representation with at most $3n+4$ positive coefficients. Equivalently,
if more masses are positive, their homogenized vectors $(v_a,1)$ are
linearly dependent. Move the masses along this dependence until one becomes
zero while all remain nonnegative, and repeat. Total mass, margins, and
the three prescribed values are preserved. This proves the support bound.
The quasi-copula formulation follows from \cref{lem:quasi}; the bivariate
facet calculation in \cref{lem:finite-bases} proves redundancy of the four
additional inequalities when $n=2$.
\end{proof}

To implement the construction for given admissible data, solve the
following finite feasibility problem, with $a\in\{1,2,3,4\}^n$:
\begin{equation}
 \begin{aligned}
 m_a&\geq0, &\sum_a m_a&=1,\\
 \sum_{a:a_i=j}m_a&=\delta_{ij}
       &&(1\leq i\leq n,\ 1\leq j\leq3),\\
 \sum_{a:a_i\leq r_{ip}\ \forall i}m_a&=q_p
       &&(1\leq p\leq3).
 \end{aligned}\label{eq:construction-lp}
\end{equation}
There are $4^n$ nonnegative unknowns and $3n+4$ equality rows. The fourth
band margins follow from total mass. A feasible basic solution gives
at most $3n+4$ positive masses. For rational data such a solution can be
chosen rational. Equations \eqref{eq:checker-formula} and
\eqref{eq:sampling} then give, respectively, the copula on the whole cube
and a sampling procedure. The cell mass $m_a$ is divided by the cell
volume to obtain its density; these are different quantities.

\section{Consequences and three-point examples}\label{sec:examples}
\subsection{A common coordinatewise upper point}
\begin{corollary}\label{cor:maximum}
If one of three points in $[0,1]^3$ is coordinatewise greater than or equal
to the other two, then quasi-copula interpolation at those points is
equivalent to copula interpolation.
\end{corollary}
\begin{proof}
Relabel the upper point as point $3$ and put it last in each coordinate
order, also when coordinates are tied. The remaining two labels can be
ordered in only two ways in each coordinate. After relabeling these two
points and permuting coordinates, the pattern is case 1 or case 3 of
\cref{tab:counts}. Neither has additional facets. Apply \cref{lem:facets}.
\end{proof}

In particular consider the ranks
\begin{equation}
 r^1=(1,2,1),\qquad r^2=(2,1,2),\qquad r^3=(3,3,3).
 \label{eq:upper-config}
\end{equation}
The coordinate orders are $(123,213,123)$, which become case 3 after
interchanging the second and third coordinates. Thus the Fr\'echet bounds
at the three points and the following directed bounds are sufficient:
\begin{align*}
 q_1-q_2&\leq\delta_{22},&
 q_2-q_1&\leq\delta_{12}+\delta_{32},\\
 q_1&\leq q_3,&q_2&\leq q_3,\\
 q_3-q_1&\leq\delta_{12}+\delta_{13}+\delta_{23}+\delta_{32}+\delta_{33},\\
 q_3-q_2&\leq\delta_{13}+\delta_{22}+\delta_{23}+\delta_{33}.
\end{align*}
There is no monotonicity requirement between $q_1$ and $q_2$.
For example the independence copula at
\[
 x^1=(0.4,0.8,0.4),\quad x^2=(0.6,0.2,0.6),\quad
 x^3=(0.9,0.9,0.9)
\]
has values $q_1=0.128>q_2=0.072$ and $q_3=0.729$.
The pattern \eqref{eq:upper-config} has $26$ facets in the reduced data
polytope. The chain pattern $(1,1,1),(2,2,2),(3,3,3)$ has $24$ and does
impose $q_1\leq q_2\leq q_3$; these are different configurations.

\subsection{Strictly interior counterexamples to pairwise sufficiency}
Take
\begin{equation}
 x^1=(0.5,0.9,0.55),\quad
 x^2=(0.55,0.5,0.9),\quad
 x^3=(0.9,0.55,0.5).\label{eq:strict-example}
\end{equation}
The orders are $(123,231,312)$, namely case 10. In each coordinate the
thresholds are $0.5,0.55,0.9$, so every point has $W_3=0$ and $M_3=0.5$.
For every ordered pair of distinct points, $D(x^a,x^b)=0.4$.
The four extra inequalities reduce to
\begin{equation}
 q_a+q_b-q_c\leq0.6\quad(\{a,b,c\}=\{1,2,3\}),
 \qquad q_1+q_2+q_3\geq0.2.\label{eq:strict-extra}
\end{equation}
Consequently $(q_1,q_2,q_3)=(0,0,0)$ admits the quasi-copula $W_3$ as an
interpolant, but violates the sum inequality. The values $(0.5,0.5,0.1)$
also satisfy every one-point and pairwise condition and hence admit the
quasi-copula \eqref{eq:quasi-extension}; they violate
$q_1+q_2-q_3\leq0.6$, since the left side is $0.9$.
Thus both types of additional constraint are necessary even when all
coordinates are strictly interior and distinct within each coordinate.

\subsection{A complete mass formula on eight cells}\label{subsec:eight}
A particularly transparent configuration is
\begin{equation}
 x^1=(1,\tfrac12,\tfrac12),\qquad
 x^2=(\tfrac12,1,\tfrac12),\qquad
 x^3=(\tfrac12,\tfrac12,1).\label{eq:half-points}
\end{equation}
The quasi-copula conditions are simply $0\leq q_p\leq\tfrac12$.
The copula criterion is
\begin{equation}
 q_1+q_2+q_3\geq\tfrac12,\qquad
 q_a+q_b-q_c\leq\tfrac12\quad(\{a,b,c\}=\{1,2,3\}).
 \label{eq:half-criterion}
\end{equation}
Its feasible region is the tetrahedron with vertices
\[
 (\tfrac12,\tfrac12,\tfrac12),\quad
 (\tfrac12,0,0),\quad(0,\tfrac12,0),\quad(0,0,\tfrac12).
\]
For instance its barycentric coefficients, in that order, are
$S-\tfrac12$, $q_1-q_2-q_3+\tfrac12$,
$q_2-q_1-q_3+\tfrac12$, and $q_3-q_1-q_2+\tfrac12$,
where $S=q_1+q_2+q_3$; they sum to one and are nonnegative exactly under
\eqref{eq:half-criterion}.

Use two bands in each coordinate, $J_{i1}=[0,\tfrac12)$ and
$J_{i2}=[\tfrac12,1)$. Let $h=m_{111}$. All masses are then forced to be
\begin{equation}
\begin{aligned}
 m_{111}&=h, &m_{112}&=q_3-h,\\
 m_{121}&=q_2-h, &m_{211}&=q_1-h,\\
 m_{122}&=\tfrac12-q_2-q_3+h,
 &m_{212}&=\tfrac12-q_1-q_3+h,\\
 m_{221}&=\tfrac12-q_1-q_2+h,
 &m_{222}&=S-\tfrac12-h.
\end{aligned}\label{eq:eight-masses}
\end{equation}
These are nonnegative precisely when
\begin{equation}
 \begin{split}
 \max\{0,q_1+q_2-\tfrac12,q_1+q_3-\tfrac12,q_2+q_3-\tfrac12\}
 \leq h\\
 \leq\min\{q_1,q_2,q_3,S-\tfrac12\}.
 \end{split}\label{eq:h-interval}
\end{equation}
Comparing the four lower bounds with the four upper bounds shows that
this interval is nonempty exactly when $0\leq q_p\leq\tfrac12$ and
\eqref{eq:half-criterion} hold. Its masses sum to one and have marginal
sums $1/2$ in each band. Moreover,
\[
 C(x^1)=m_{111}+m_{211}=q_1,\quad
 C(x^2)=m_{111}+m_{121}=q_2,\quad
 C(x^3)=m_{111}+m_{112}=q_3.
\]
Thus \eqref{eq:eight-masses}, with any $h$ in \eqref{eq:h-interval},
and \eqref{eq:checker-formula} explicitly construct the interpolant.

For equal prescribed values $q_1=q_2=q_3=t$, feasibility is exactly
$1/6\leq t\leq1/2$. At $t=1/6$ the only possible $h$ is zero:
put mass $1/6$ in each of the six cells other than $111$ and $222$,
and zero in those two cells. Each cell has volume $1/8$, so the copula
density is $4/3$ on the six occupied cells and zero on the other two.
This illustrates the construction without solving a linear program.

\clearpage
\appendix
\raggedbottom
% BEGIN GENERATED APPENDIX TABLES
\section{Order configurations and additional facets in dimensions two to six}
\label{app:all-tables}
This appendix lists every symmetry type in dimensions $2\leq n\leq6$.
Two order lists represent the same type when they differ by a permutation
of coordinates and a common relabeling of the three points. We choose the
lexicographically least sorted representative. Case numbers start anew in
each dimension. The trivariate rows repeat Tables~\ref{tab:counts}
and~\ref{tab:extras}, so that all five dimensions can be read together.

In each coordinate write
\[
0=t_{i0}\leq t_{i1}\leq t_{i2}\leq t_{i3}\leq t_{i4}=1,
\qquad d_i=t_{i2}-t_{i1}.
\]
An order string $abc$ means
$x_i^a=t_{i1}$, $x_i^b=t_{i2}$, and $x_i^c=t_{i3}$.
As in the trivariate tables, put
\[
U_{ab\mid c}=q_a+q_b-q_c,\qquad
S=q_1+q_2+q_3,\qquad
B_n=\sum_{i=1}^n(t_{i1}+t_{i3})-(2n-2).
\]
The \emph{baseline} consists of the grid inequalities, the one-point
Fr\'echet--Hoeffding bounds, and the directed pairwise bounds. The count
tables give the number of distinct facets of each variable-data polytope,
not the number of inequalities remaining after a particular grid is fixed.
An orbit size counts labeled coordinate-order lists; its sum is $6^n$.

The companion table lists precisely the facets outside the baseline.
An expression $U\leq\min\{f_1,\ldots,f_k\}+g$ represents $k$ separate
linear inequalities $U\leq f_j+g$. Each listed constituent is a facet.
An entry of \emph{None} means that the baseline already describes the
entire polytope. Thus the baseline and the displayed additional inequalities
form a complete description for the given order configuration.
The lower-bound constants $2n-2$ are $2,4,6,8,10$ in the five dimensions.

\begin{table}[htbp]
\centering\small
\caption{Totals over all representatives in dimensions two to six.}
\label{tab:all-dimension-totals}
\begin{tabular}{rrrrrrr}
\toprule
$n$ & Order lists & Types & Facets & Baseline & Upper & Lower\\
\midrule
2 & 36 & 5 & 83 & 83 & 0 & 0\\
3 & 216 & 10 & 283 & 263 & 18 & 2\\
4 & 1296 & 24 & 913 & 802 & 104 & 7\\
5 & 7776 & 42 & 2004 & 1682 & 305 & 17\\
6 & 46656 & 83 & 4781 & 3865 & 874 & 42\\
\bottomrule
\end{tabular}
\end{table}
The Upper and Lower columns count additional facets of the corresponding
three-point families; their sum is the total outside the baseline.
All totals are sums over one representative per symmetry type.

\clearpage
\subsection{Dimension 2}\label{app:tables-n2}
There are $ 5$ symmetry types, covering all $ 36$ order lists.
{\small\setlength{\tabcolsep}{4.5pt}
\begin{longtable}{@{}rlrrrr@{}}
\caption{Dimension $ 2$: order configurations and exact facet counts.}\label{tab:counts-n2}\\
\toprule
Case & Coordinate orders & Orbit & Facets & Baseline & Additional\\
\midrule\endfirsthead
\multicolumn{6}{l}{\small Table \ref{tab:counts-n2} (continued).}\\
\toprule
Case & Coordinate orders & Orbit & Facets & Baseline & Additional\\
\midrule\endhead
\midrule
\multicolumn{6}{r}{\small Continued on next page.}\\
\endfoot\bottomrule\endlastfoot
1 & $(123,123)$ & 6 & 16 & 16 & 0\\
2 & $(123,132)$ & 6 & 17 & 17 & 0\\
3 & $(123,213)$ & 6 & 17 & 17 & 0\\
4 & $(123,231)$ & 12 & 17 & 17 & 0\\
5 & $(123,321)$ & 6 & 16 & 16 & 0\\
\midrule
 & Total & 36 & 83 & 83 & 0\\
\end{longtable}
}

{\small\setlength{\tabcolsep}{4.5pt}
\begin{longtable}{@{}r>{\raggedright\arraybackslash}p{12.9cm}@{}}
\caption{Dimension $ 2$: all additional facet inequalities beyond the baseline.}\label{tab:extras-n2}\\
\toprule
Case & Additional inequalities\\
\midrule\endfirsthead
\multicolumn{2}{l}{\small Table \ref{tab:extras-n2} (continued).}\\
\toprule
Case & Additional inequalities\\
\midrule\endhead
\midrule
\multicolumn{2}{r}{\small Continued on next page.}\\
\endfoot\bottomrule\endlastfoot
1 & None.\\[3pt]
2 & None.\\[3pt]
3 & None.\\[3pt]
4 & None.\\[3pt]
5 & None.\\[3pt]
\end{longtable}
}

\clearpage
\subsection{Dimension 3}\label{app:tables-n3}
There are $ 10$ symmetry types, covering all $ 216$ order lists.
{\small\setlength{\tabcolsep}{4.5pt}
\begin{longtable}{@{}rlrrrr@{}}
\caption{Dimension $ 3$: order configurations and exact facet counts.}\label{tab:counts-n3}\\
\toprule
Case & Coordinate orders & Orbit & Facets & Baseline & Additional\\
\midrule\endfirsthead
\multicolumn{6}{l}{\small Table \ref{tab:counts-n3} (continued).}\\
\toprule
Case & Coordinate orders & Orbit & Facets & Baseline & Additional\\
\midrule\endhead
\midrule
\multicolumn{6}{r}{\small Continued on next page.}\\
\endfoot\bottomrule\endlastfoot
1 & $(123,123,123)$ & 6 & 24 & 24 & 0\\
2 & $(123,123,132)$ & 18 & 28 & 26 & 2\\
3 & $(123,123,213)$ & 18 & 26 & 26 & 0\\
4 & $(123,123,231)$ & 18 & 29 & 27 & 2\\
5 & $(123,123,312)$ & 18 & 28 & 26 & 2\\
6 & $(123,123,321)$ & 18 & 28 & 26 & 2\\
7 & $(123,132,213)$ & 36 & 28 & 26 & 2\\
8 & $(123,132,231)$ & 36 & 32 & 28 & 4\\
9 & $(123,213,312)$ & 36 & 29 & 27 & 2\\
10 & $(123,231,312)$ & 12 & 31 & 27 & 4\\
\midrule
 & Total & 216 & 283 & 263 & 20\\
\end{longtable}
}

{\small\setlength{\tabcolsep}{4.5pt}
\begin{longtable}{@{}r>{\raggedright\arraybackslash}p{12.9cm}@{}}
\caption{Dimension $ 3$: all additional facet inequalities beyond the baseline.}\label{tab:extras-n3}\\
\toprule
Case & Additional inequalities\\
\midrule\endfirsthead
\multicolumn{2}{l}{\small Table \ref{tab:extras-n3} (continued).}\\
\toprule
Case & Additional inequalities\\
\midrule\endhead
\midrule
\multicolumn{2}{r}{\small Continued on next page.}\\
\endfoot\bottomrule\endlastfoot
1 & None.\\[3pt]
2 & \(U_{23\mid 1}\leq \min\{t_{13},t_{23}\}+d_{1}+d_{2}+d_{3}\)\\[3pt]
3 & None.\\[3pt]
4 & \(U_{13\mid 2}\leq \min\{t_{13},t_{23}\}+d_{3}\)\\[3pt]
5 & \(U_{23\mid 1}\leq \min\{t_{13},t_{23}\}+d_{1}+d_{2}\)\\[3pt]
6 & \(U_{13\mid 2}\leq \min\{t_{13},t_{23}\}\)\\[3pt]
7 & \(U_{23\mid 1}\leq \min\{t_{13},t_{33}\}+d_{1}+d_{2}\)\\[3pt]
8 & \(U_{23\mid 1}\leq t_{32}+d_{1}+d_{2}\)\newline \(U_{13\mid 2}\leq t_{22}+d_{3}\)\newline \(U_{12\mid 3}\leq t_{12}\)\newline \(S\geq B_n\)\\[3pt]
9 & \(U_{23\mid 1}\leq \min\{t_{13},t_{23}\}+d_{1}\)\\[3pt]
10 & \(U_{23\mid 1}\leq t_{22}+d_{1}\)\newline \(U_{13\mid 2}\leq t_{32}+d_{2}\)\newline \(U_{12\mid 3}\leq t_{12}+d_{3}\)\newline \(S\geq B_n\)\\[3pt]
\end{longtable}
}

\clearpage
\subsection{Dimension 4}\label{app:tables-n4}
There are $ 24$ symmetry types, covering all $ 1296$ order lists.
{\small\setlength{\tabcolsep}{4.5pt}
\begin{longtable}{@{}rlrrrr@{}}
\caption{Dimension $ 4$: order configurations and exact facet counts.}\label{tab:counts-n4}\\
\toprule
Case & Coordinate orders & Orbit & Facets & Baseline & Additional\\
\midrule\endfirsthead
\multicolumn{6}{l}{\small Table \ref{tab:counts-n4} (continued).}\\
\toprule
Case & Coordinate orders & Orbit & Facets & Baseline & Additional\\
\midrule\endhead
\midrule
\multicolumn{6}{r}{\small Continued on next page.}\\
\endfoot\bottomrule\endlastfoot
1 & $(123,123,123,123)$ & 6 & 30 & 30 & 0\\
2 & $(123,123,123,132)$ & 24 & 35 & 32 & 3\\
3 & $(123,123,123,213)$ & 24 & 32 & 32 & 0\\
4 & $(123,123,123,231)$ & 24 & 36 & 33 & 3\\
5 & $(123,123,123,312)$ & 24 & 35 & 32 & 3\\
6 & $(123,123,123,321)$ & 24 & 35 & 32 & 3\\
7 & $(123,123,132,132)$ & 18 & 37 & 33 & 4\\
8 & $(123,123,132,213)$ & 72 & 35 & 32 & 3\\
9 & $(123,123,132,231)$ & 72 & 43 & 34 & 9\\
10 & $(123,123,132,312)$ & 72 & 37 & 33 & 4\\
11 & $(123,123,132,321)$ & 72 & 44 & 35 & 9\\
12 & $(123,123,213,213)$ & 18 & 33 & 33 & 0\\
13 & $(123,123,213,231)$ & 72 & 36 & 33 & 3\\
14 & $(123,123,213,312)$ & 72 & 36 & 33 & 3\\
15 & $(123,123,213,321)$ & 72 & 37 & 34 & 3\\
16 & $(123,123,231,231)$ & 36 & 38 & 34 & 4\\
17 & $(123,123,231,312)$ & 72 & 43 & 34 & 9\\
18 & $(123,123,231,321)$ & 72 & 39 & 35 & 4\\
19 & $(123,123,312,321)$ & 72 & 44 & 35 & 9\\
20 & $(123,123,321,321)$ & 18 & 38 & 34 & 4\\
21 & $(123,132,213,231)$ & 72 & 43 & 34 & 9\\
22 & $(123,132,213,312)$ & 72 & 38 & 34 & 4\\
23 & $(123,132,213,321)$ & 144 & 44 & 35 & 9\\
24 & $(123,132,231,321)$ & 72 & 45 & 36 & 9\\
\midrule
 & Total & 1296 & 913 & 802 & 111\\
\end{longtable}
}

\clearpage
{\small\setlength{\tabcolsep}{4.5pt}
\begin{longtable}{@{}r>{\raggedright\arraybackslash}p{12.9cm}@{}}
\caption{Dimension $ 4$: all additional facet inequalities beyond the baseline.}\label{tab:extras-n4}\\
\toprule
Case & Additional inequalities\\
\midrule\endfirsthead
\multicolumn{2}{l}{\small Table \ref{tab:extras-n4} (continued).}\\
\toprule
Case & Additional inequalities\\
\midrule\endhead
\midrule
\multicolumn{2}{r}{\small Continued on next page.}\\
\endfoot\bottomrule\endlastfoot
1 & None.\\
2 & \(U_{23\mid 1}\leq \min\{t_{13},t_{23},t_{33}\}+d_{1}+d_{2}+d_{3}+d_{4}\)\\
3 & None.\\
4 & \(U_{13\mid 2}\leq \min\{t_{13},t_{23},t_{33}\}+d_{4}\)\\
5 & \(U_{23\mid 1}\leq \min\{t_{13},t_{23},t_{33}\}+d_{1}+d_{2}+d_{3}\)\\
6 & \(U_{13\mid 2}\leq \min\{t_{13},t_{23},t_{33}\}\)\\
7 & \(U_{23\mid 1}\leq \min\{t_{13},t_{23},t_{33},t_{43}\}+d_{1}+d_{2}+d_{3}+d_{4}\)\\
8 & \(U_{23\mid 1}\leq \min\{t_{13},t_{23},t_{43}\}+d_{1}+d_{2}+d_{3}\)\\
9 & \(U_{23\mid 1}\leq \min\{t_{13},t_{23},t_{42}\}+d_{1}+d_{2}+d_{3}\)\newline \(U_{13\mid 2}\leq \min\{t_{13},t_{23},t_{32}\}+d_{4}\)\newline \(U_{12\mid 3}\leq \min\{t_{12},t_{22}\}\)\newline \(S\geq B_n\)\\
10 & \(U_{23\mid 1}\leq \min\{t_{13},t_{23},t_{33},t_{43}\}+d_{1}+d_{2}+d_{3}\)\\
11 & \(U_{23\mid 1}\leq \min\{t_{13},t_{23},t_{42}\}+d_{1}+d_{2}+d_{3}\)\newline \(U_{13\mid 2}\leq \min\{t_{13},t_{23},t_{32}\}\)\newline \(U_{12\mid 3}\leq \min\{t_{12},t_{22}\}+d_{4}\)\newline \(S\geq B_n\)\\
12 & None.\\
13 & \(U_{13\mid 2}\leq \min\{t_{13},t_{23},t_{33}\}+d_{3}+d_{4}\)\\
14 & \(U_{23\mid 1}\leq \min\{t_{13},t_{23},t_{33}\}+d_{1}+d_{2}\)\\
15 & \(U_{13\mid 2}\leq \min\{t_{13},t_{23},t_{33}\}+d_{3}\)\\
16 & \(U_{13\mid 2}\leq \min\{t_{13},t_{23},t_{33},t_{43}\}+d_{3}+d_{4}\)\\
17 & \(U_{23\mid 1}\leq \min\{t_{13},t_{23},t_{32}\}+d_{1}+d_{2}\)\newline \(U_{13\mid 2}\leq \min\{t_{13},t_{23},t_{42}\}+d_{3}\)\newline \(U_{12\mid 3}\leq \min\{t_{12},t_{22}\}+d_{4}\)\newline \(S\geq B_n\)\\
18 & \(U_{13\mid 2}\leq \min\{t_{13},t_{23},t_{33},t_{43}\}+d_{3}\)\\
19 & \(U_{23\mid 1}\leq \min\{t_{13},t_{23},t_{42}\}+d_{1}+d_{2}\)\newline \(U_{13\mid 2}\leq \min\{t_{13},t_{23},t_{32}\}\)\newline \(U_{12\mid 3}\leq \min\{t_{12},t_{22}\}+d_{3}+d_{4}\)\newline \(S\geq B_n\)\\
20 & \(U_{13\mid 2}\leq \min\{t_{13},t_{23},t_{33},t_{43}\}\)\\
21 & \(U_{23\mid 1}\leq \min\{t_{13},t_{33},t_{42}\}+d_{1}+d_{2}\)\newline \(U_{13\mid 2}\leq \min\{t_{13},t_{22},t_{33}\}+d_{3}+d_{4}\)\newline \(U_{12\mid 3}\leq \min\{t_{12},t_{32}\}\)\newline \(S\geq B_n\)\\
22 & \(U_{23\mid 1}\leq \min\{t_{13},t_{23},t_{33},t_{43}\}+d_{1}+d_{2}\)\\
23 & \(U_{23\mid 1}\leq \min\{t_{13},t_{33},t_{42}\}+d_{1}+d_{2}\)\newline \(U_{13\mid 2}\leq \min\{t_{13},t_{22},t_{33}\}+d_{3}\)\newline \(U_{12\mid 3}\leq \min\{t_{12},t_{32}\}+d_{4}\)\newline \(S\geq B_n\)\\
24 & \(U_{23\mid 1}\leq \min\{t_{32},t_{42}\}+d_{1}+d_{2}\)\newline \(U_{13\mid 2}\leq \min\{t_{22},t_{33},t_{43}\}+d_{3}\)\newline \(U_{12\mid 3}\leq \min\{t_{12},t_{33},t_{43}\}+d_{4}\)\newline \(S\geq B_n\)\\
\end{longtable}
}

\clearpage
\subsection{Dimension 5}\label{app:tables-n5}
There are $ 42$ symmetry types, covering all $ 7776$ order lists.
{\small\setlength{\tabcolsep}{4.5pt}
\begin{longtable}{@{}rlrrrr@{}}
\caption{Dimension $ 5$: order configurations and exact facet counts.}\label{tab:counts-n5}\\
\toprule
Case & Coordinate orders & Orbit & Facets & Baseline & Additional\\
\midrule\endfirsthead
\multicolumn{6}{l}{\small Table \ref{tab:counts-n5} (continued).}\\
\toprule
Case & Coordinate orders & Orbit & Facets & Baseline & Additional\\
\midrule\endhead
\midrule
\multicolumn{6}{r}{\small Continued on next page.}\\
\endfoot\bottomrule\endlastfoot
1 & $(123,123,123,123,123)$ & 6 & 36 & 36 & 0\\
2 & $(123,123,123,123,132)$ & 30 & 42 & 38 & 4\\
3 & $(123,123,123,123,213)$ & 30 & 38 & 38 & 0\\
4 & $(123,123,123,123,231)$ & 30 & 43 & 39 & 4\\
5 & $(123,123,123,123,312)$ & 30 & 42 & 38 & 4\\
6 & $(123,123,123,123,321)$ & 30 & 42 & 38 & 4\\
7 & $(123,123,123,132,132)$ & 60 & 44 & 39 & 5\\
8 & $(123,123,123,132,213)$ & 120 & 42 & 38 & 4\\
9 & $(123,123,123,132,231)$ & 120 & 52 & 40 & 12\\
10 & $(123,123,123,132,312)$ & 120 & 44 & 39 & 5\\
11 & $(123,123,123,132,321)$ & 120 & 53 & 41 & 12\\
12 & $(123,123,123,213,213)$ & 60 & 39 & 39 & 0\\
13 & $(123,123,123,213,231)$ & 120 & 43 & 39 & 4\\
14 & $(123,123,123,213,312)$ & 120 & 43 & 39 & 4\\
15 & $(123,123,123,213,321)$ & 120 & 44 & 40 & 4\\
16 & $(123,123,123,231,231)$ & 60 & 45 & 40 & 5\\
17 & $(123,123,123,231,312)$ & 120 & 52 & 40 & 12\\
18 & $(123,123,123,231,321)$ & 120 & 46 & 41 & 5\\
19 & $(123,123,123,312,312)$ & 60 & 45 & 40 & 5\\
20 & $(123,123,123,312,321)$ & 120 & 53 & 41 & 12\\
21 & $(123,123,123,321,321)$ & 60 & 45 & 40 & 5\\
22 & $(123,123,132,132,213)$ & 180 & 44 & 39 & 5\\
23 & $(123,123,132,132,231)$ & 180 & 55 & 41 & 14\\
24 & $(123,123,132,213,213)$ & 180 & 43 & 39 & 4\\
25 & $(123,123,132,213,231)$ & 360 & 52 & 40 & 12\\
26 & $(123,123,132,213,312)$ & 360 & 45 & 40 & 5\\
27 & $(123,123,132,213,321)$ & 360 & 53 & 41 & 12\\
28 & $(123,123,132,231,231)$ & 180 & 55 & 41 & 14\\
29 & $(123,123,132,231,312)$ & 360 & 55 & 41 & 14\\
30 & $(123,123,132,231,321)$ & 360 & 56 & 42 & 14\\
31 & $(123,123,132,312,312)$ & 180 & 45 & 40 & 5\\
32 & $(123,123,132,312,321)$ & 360 & 55 & 41 & 14\\
33 & $(123,123,132,321,321)$ & 180 & 56 & 42 & 14\\
34 & $(123,123,213,213,312)$ & 180 & 44 & 40 & 4\\
35 & $(123,123,213,231,312)$ & 360 & 53 & 41 & 12\\
36 & $(123,123,213,231,321)$ & 360 & 46 & 41 & 5\\
37 & $(123,123,213,312,312)$ & 180 & 46 & 41 & 5\\
38 & $(123,123,213,312,321)$ & 360 & 54 & 42 & 12\\
39 & $(123,123,213,321,321)$ & 180 & 47 & 42 & 5\\
40 & $(123,123,231,231,312)$ & 180 & 55 & 41 & 14\\
41 & $(123,123,231,312,321)$ & 360 & 56 & 42 & 14\\
42 & $(123,132,213,231,312)$ & 720 & 56 & 42 & 14\\
\midrule
 & Total & 7776 & 2004 & 1682 & 322\\
\end{longtable}
}

\clearpage
{\small\setlength{\tabcolsep}{4.5pt}
\begin{longtable}{@{}r>{\raggedright\arraybackslash}p{12.9cm}@{}}
\caption{Dimension $ 5$: all additional facet inequalities beyond the baseline.}\label{tab:extras-n5}\\
\toprule
Case & Additional inequalities\\
\midrule\endfirsthead
\multicolumn{2}{l}{\small Table \ref{tab:extras-n5} (continued).}\\
\toprule
Case & Additional inequalities\\
\midrule\endhead
\midrule
\multicolumn{2}{r}{\small Continued on next page.}\\
\endfoot\bottomrule\endlastfoot
1 & None.\\[3pt]
2 & \(U_{23\mid 1}\leq \min\{t_{13},t_{23},t_{33},t_{43}\}+d_{1}+d_{2}+d_{3}+d_{4}+d_{5}\)\\[3pt]
3 & None.\\[3pt]
4 & \(U_{13\mid 2}\leq \min\{t_{13},t_{23},t_{33},t_{43}\}+d_{5}\)\\[3pt]
5 & \(U_{23\mid 1}\leq \min\{t_{13},t_{23},t_{33},t_{43}\}+d_{1}+d_{2}+d_{3}+d_{4}\)\\[3pt]
6 & \(U_{13\mid 2}\leq \min\{t_{13},t_{23},t_{33},t_{43}\}\)\\[3pt]
7 & \(U_{23\mid 1}\leq \min\{t_{13},t_{23},t_{33},t_{43},t_{53}\}+d_{1}+d_{2}+d_{3}+d_{4}+d_{5}\)\\[3pt]
8 & \(U_{23\mid 1}\leq \min\{t_{13},t_{23},t_{33},t_{53}\}+d_{1}+d_{2}+d_{3}+d_{4}\)\\[3pt]
9 & \(U_{23\mid 1}\leq \min\{t_{13},t_{23},t_{33},t_{52}\}+d_{1}+d_{2}+d_{3}+d_{4}\)\newline \(U_{13\mid 2}\leq \min\{t_{13},t_{23},t_{33},t_{42}\}+d_{5}\)\newline \(U_{12\mid 3}\leq \min\{t_{12},t_{22},t_{32}\}\)\newline \(S\geq B_n\)\\[3pt]
10 & \(U_{23\mid 1}\leq \min\{t_{13},t_{23},t_{33},t_{43},t_{53}\}+d_{1}+d_{2}+d_{3}+d_{4}\)\\[3pt]
11 & \(U_{23\mid 1}\leq \min\{t_{13},t_{23},t_{33},t_{52}\}+d_{1}+d_{2}+d_{3}+d_{4}\)\newline \(U_{13\mid 2}\leq \min\{t_{13},t_{23},t_{33},t_{42}\}\)\newline \(U_{12\mid 3}\leq \min\{t_{12},t_{22},t_{32}\}+d_{5}\)\newline \(S\geq B_n\)\\[3pt]
12 & None.\\[3pt]
13 & \(U_{13\mid 2}\leq \min\{t_{13},t_{23},t_{33},t_{43}\}+d_{4}+d_{5}\)\\[3pt]
14 & \(U_{23\mid 1}\leq \min\{t_{13},t_{23},t_{33},t_{43}\}+d_{1}+d_{2}+d_{3}\)\\[3pt]
15 & \(U_{13\mid 2}\leq \min\{t_{13},t_{23},t_{33},t_{43}\}+d_{4}\)\\[3pt]
16 & \(U_{13\mid 2}\leq \min\{t_{13},t_{23},t_{33},t_{43},t_{53}\}+d_{4}+d_{5}\)\\[3pt]
17 & \(U_{23\mid 1}\leq \min\{t_{13},t_{23},t_{33},t_{42}\}+d_{1}+d_{2}+d_{3}\)\newline \(U_{13\mid 2}\leq \min\{t_{13},t_{23},t_{33},t_{52}\}+d_{4}\)\newline \(U_{12\mid 3}\leq \min\{t_{12},t_{22},t_{32}\}+d_{5}\)\newline \(S\geq B_n\)\\[3pt]
18 & \(U_{13\mid 2}\leq \min\{t_{13},t_{23},t_{33},t_{43},t_{53}\}+d_{4}\)\\[3pt]
19 & \(U_{23\mid 1}\leq \min\{t_{13},t_{23},t_{33},t_{43},t_{53}\}+d_{1}+d_{2}+d_{3}\)\\[3pt]
20 & \(U_{23\mid 1}\leq \min\{t_{13},t_{23},t_{33},t_{52}\}+d_{1}+d_{2}+d_{3}\)\newline \(U_{13\mid 2}\leq \min\{t_{13},t_{23},t_{33},t_{42}\}\)\newline \(U_{12\mid 3}\leq \min\{t_{12},t_{22},t_{32}\}+d_{4}+d_{5}\)\newline \(S\geq B_n\)\\[3pt]
21 & \(U_{13\mid 2}\leq \min\{t_{13},t_{23},t_{33},t_{43},t_{53}\}\)\\[3pt]
22 & \(U_{23\mid 1}\leq \min\{t_{13},t_{23},t_{33},t_{43},t_{53}\}+d_{1}+d_{2}+d_{3}+d_{4}\)\\[3pt]
23 & \(U_{23\mid 1}\leq \min\{t_{13},t_{23},t_{33},t_{43},t_{52}\}+d_{1}+d_{2}+d_{3}+d_{4}\)\newline \(U_{13\mid 2}\leq \min\{t_{13},t_{23},t_{32},t_{42}\}+d_{5}\)\newline \(U_{12\mid 3}\leq \min\{t_{12},t_{22},t_{33},t_{43}\}\)\newline \(S\geq B_n\)\\[3pt]
24 & \(U_{23\mid 1}\leq \min\{t_{13},t_{23},t_{43},t_{53}\}+d_{1}+d_{2}+d_{3}\)\\[3pt]
25 & \(U_{23\mid 1}\leq \min\{t_{13},t_{23},t_{43},t_{52}\}+d_{1}+d_{2}+d_{3}\)\newline \(U_{13\mid 2}\leq \min\{t_{13},t_{23},t_{32},t_{43}\}+d_{4}+d_{5}\)\newline \(U_{12\mid 3}\leq \min\{t_{12},t_{22},t_{42}\}\)\newline \(S\geq B_n\)\\[3pt]
26 & \(U_{23\mid 1}\leq \min\{t_{13},t_{23},t_{33},t_{43},t_{53}\}+d_{1}+d_{2}+d_{3}\)\\[3pt]
27 & \(U_{23\mid 1}\leq \min\{t_{13},t_{23},t_{43},t_{52}\}+d_{1}+d_{2}+d_{3}\)\newline \(U_{13\mid 2}\leq \min\{t_{13},t_{23},t_{32},t_{43}\}+d_{4}\)\newline \(U_{12\mid 3}\leq \min\{t_{12},t_{22},t_{42}\}+d_{5}\)\newline \(S\geq B_n\)\\[3pt]
28 & \(U_{23\mid 1}\leq \min\{t_{13},t_{23},t_{42},t_{52}\}+d_{1}+d_{2}+d_{3}\)\newline \(U_{13\mid 2}\leq \min\{t_{13},t_{23},t_{32},t_{43},t_{53}\}+d_{4}+d_{5}\)\newline \(U_{12\mid 3}\leq \min\{t_{12},t_{22},t_{43},t_{53}\}\)\newline \(S\geq B_n\)\\[3pt]
29 & \(U_{23\mid 1}\leq \min\{t_{13},t_{23},t_{33},t_{42},t_{53}\}+d_{1}+d_{2}+d_{3}\)\newline \(U_{13\mid 2}\leq \min\{t_{13},t_{23},t_{32},t_{52}\}+d_{4}\)\newline \(U_{12\mid 3}\leq \min\{t_{12},t_{22},t_{33},t_{53}\}+d_{5}\)\newline \(S\geq B_n\)\\[3pt]
30 & \(U_{23\mid 1}\leq \min\{t_{13},t_{23},t_{42},t_{52}\}+d_{1}+d_{2}+d_{3}\)\newline \(U_{13\mid 2}\leq \min\{t_{13},t_{23},t_{32},t_{43},t_{53}\}+d_{4}\)\newline \(U_{12\mid 3}\leq \min\{t_{12},t_{22},t_{43},t_{53}\}+d_{5}\)\newline \(S\geq B_n\)\\[3pt]
31 & \(U_{23\mid 1}\leq \min\{t_{13},t_{23},t_{33},t_{43},t_{53}\}+d_{1}+d_{2}+d_{3}\)\\[3pt]
32 & \(U_{23\mid 1}\leq \min\{t_{13},t_{23},t_{33},t_{43},t_{52}\}+d_{1}+d_{2}+d_{3}\)\newline \(U_{13\mid 2}\leq \min\{t_{13},t_{23},t_{32},t_{42}\}\)\newline \(U_{12\mid 3}\leq \min\{t_{12},t_{22},t_{33},t_{43}\}+d_{4}+d_{5}\)\newline \(S\geq B_n\)\\[3pt]
33 & \(U_{23\mid 1}\leq \min\{t_{13},t_{23},t_{42},t_{52}\}+d_{1}+d_{2}+d_{3}\)\newline \(U_{13\mid 2}\leq \min\{t_{13},t_{23},t_{32},t_{43},t_{53}\}\)\newline \(U_{12\mid 3}\leq \min\{t_{12},t_{22},t_{43},t_{53}\}+d_{4}+d_{5}\)\newline \(S\geq B_n\)\\[3pt]
34 & \(U_{23\mid 1}\leq \min\{t_{13},t_{23},t_{33},t_{43}\}+d_{1}+d_{2}\)\\[3pt]
35 & \(U_{23\mid 1}\leq \min\{t_{13},t_{23},t_{33},t_{42}\}+d_{1}+d_{2}\)\newline \(U_{13\mid 2}\leq \min\{t_{13},t_{23},t_{33},t_{52}\}+d_{3}+d_{4}\)\newline \(U_{12\mid 3}\leq \min\{t_{12},t_{22},t_{32}\}+d_{5}\)\newline \(S\geq B_n\)\\[3pt]
36 & \(U_{13\mid 2}\leq \min\{t_{13},t_{23},t_{33},t_{43},t_{53}\}+d_{3}+d_{4}\)\\[3pt]
37 & \(U_{23\mid 1}\leq \min\{t_{13},t_{23},t_{33},t_{43},t_{53}\}+d_{1}+d_{2}\)\\[3pt]
38 & \(U_{23\mid 1}\leq \min\{t_{13},t_{23},t_{33},t_{52}\}+d_{1}+d_{2}\)\newline \(U_{13\mid 2}\leq \min\{t_{13},t_{23},t_{33},t_{42}\}+d_{3}\)\newline \(U_{12\mid 3}\leq \min\{t_{12},t_{22},t_{32}\}+d_{4}+d_{5}\)\newline \(S\geq B_n\)\\[3pt]
39 & \(U_{13\mid 2}\leq \min\{t_{13},t_{23},t_{33},t_{43},t_{53}\}+d_{3}\)\\[3pt]
40 & \(U_{23\mid 1}\leq \min\{t_{13},t_{23},t_{32},t_{42}\}+d_{1}+d_{2}\)\newline \(U_{13\mid 2}\leq \min\{t_{13},t_{23},t_{33},t_{43},t_{52}\}+d_{3}+d_{4}\)\newline \(U_{12\mid 3}\leq \min\{t_{12},t_{22},t_{33},t_{43}\}+d_{5}\)\newline \(S\geq B_n\)\\[3pt]
41 & \(U_{23\mid 1}\leq \min\{t_{13},t_{23},t_{32},t_{52}\}+d_{1}+d_{2}\)\newline \(U_{13\mid 2}\leq \min\{t_{13},t_{23},t_{33},t_{42},t_{53}\}+d_{3}\)\newline \(U_{12\mid 3}\leq \min\{t_{12},t_{22},t_{33},t_{53}\}+d_{4}+d_{5}\)\newline \(S\geq B_n\)\\[3pt]
42 & \(U_{23\mid 1}\leq \min\{t_{13},t_{23},t_{33},t_{42},t_{53}\}+d_{1}+d_{2}\)\newline \(U_{13\mid 2}\leq \min\{t_{13},t_{22},t_{33},t_{52}\}+d_{3}+d_{4}\)\newline \(U_{12\mid 3}\leq \min\{t_{12},t_{23},t_{32},t_{53}\}+d_{5}\)\newline \(S\geq B_n\)\\[3pt]
\end{longtable}
}

\clearpage
\subsection{Dimension 6}\label{app:tables-n6}
There are $ 83$ symmetry types, covering all $ 46656$ order lists.
{\small\setlength{\tabcolsep}{4.5pt}
\begin{longtable}{@{}rlrrrr@{}}
\caption{Dimension $ 6$: order configurations and exact facet counts.}\label{tab:counts-n6}\\
\toprule
Case & Coordinate orders & Orbit & Facets & Baseline & Additional\\
\midrule\endfirsthead
\multicolumn{6}{l}{\small Table \ref{tab:counts-n6} (continued).}\\
\toprule
Case & Coordinate orders & Orbit & Facets & Baseline & Additional\\
\midrule\endhead
\midrule
\multicolumn{6}{r}{\small Continued on next page.}\\
\endfoot\bottomrule\endlastfoot
1 & $(123,123,123,123,123,123)$ & 6 & 42 & 42 & 0\\
2 & $(123,123,123,123,123,132)$ & 36 & 49 & 44 & 5\\
3 & $(123,123,123,123,123,213)$ & 36 & 44 & 44 & 0\\
4 & $(123,123,123,123,123,231)$ & 36 & 50 & 45 & 5\\
5 & $(123,123,123,123,123,312)$ & 36 & 49 & 44 & 5\\
6 & $(123,123,123,123,123,321)$ & 36 & 49 & 44 & 5\\
7 & $(123,123,123,123,132,132)$ & 90 & 51 & 45 & 6\\
8 & $(123,123,123,123,132,213)$ & 180 & 49 & 44 & 5\\
9 & $(123,123,123,123,132,231)$ & 180 & 61 & 46 & 15\\
10 & $(123,123,123,123,132,312)$ & 180 & 51 & 45 & 6\\
11 & $(123,123,123,123,132,321)$ & 180 & 62 & 47 & 15\\
12 & $(123,123,123,123,213,213)$ & 90 & 45 & 45 & 0\\
13 & $(123,123,123,123,213,231)$ & 180 & 50 & 45 & 5\\
14 & $(123,123,123,123,213,312)$ & 180 & 50 & 45 & 5\\
15 & $(123,123,123,123,213,321)$ & 180 & 51 & 46 & 5\\
16 & $(123,123,123,123,231,231)$ & 90 & 52 & 46 & 6\\
17 & $(123,123,123,123,231,312)$ & 180 & 61 & 46 & 15\\
18 & $(123,123,123,123,231,321)$ & 180 & 53 & 47 & 6\\
19 & $(123,123,123,123,312,312)$ & 90 & 52 & 46 & 6\\
20 & $(123,123,123,123,312,321)$ & 180 & 62 & 47 & 15\\
21 & $(123,123,123,123,321,321)$ & 90 & 52 & 46 & 6\\
22 & $(123,123,123,132,132,132)$ & 60 & 51 & 45 & 6\\
23 & $(123,123,123,132,132,213)$ & 360 & 51 & 45 & 6\\
24 & $(123,123,123,132,132,231)$ & 360 & 64 & 47 & 17\\
25 & $(123,123,123,132,132,312)$ & 360 & 51 & 45 & 6\\
26 & $(123,123,123,132,132,321)$ & 360 & 64 & 47 & 17\\
27 & $(123,123,123,132,213,213)$ & 360 & 50 & 45 & 5\\
28 & $(123,123,123,132,213,231)$ & 720 & 61 & 46 & 15\\
29 & $(123,123,123,132,213,312)$ & 720 & 52 & 46 & 6\\
30 & $(123,123,123,132,213,321)$ & 720 & 62 & 47 & 15\\
31 & $(123,123,123,132,231,231)$ & 360 & 64 & 47 & 17\\
32 & $(123,123,123,132,231,312)$ & 720 & 64 & 47 & 17\\
33 & $(123,123,123,132,231,321)$ & 720 & 65 & 48 & 17\\
34 & $(123,123,123,132,312,312)$ & 360 & 52 & 46 & 6\\
35 & $(123,123,123,132,312,321)$ & 720 & 64 & 47 & 17\\
36 & $(123,123,123,132,321,321)$ & 360 & 65 & 48 & 17\\
37 & $(123,123,123,213,213,213)$ & 60 & 45 & 45 & 0\\
38 & $(123,123,123,213,213,231)$ & 360 & 50 & 45 & 5\\
39 & $(123,123,123,213,213,312)$ & 360 & 51 & 46 & 5\\
40 & $(123,123,123,213,213,321)$ & 360 & 51 & 46 & 5\\
41 & $(123,123,123,213,231,231)$ & 360 & 52 & 46 & 6\\
42 & $(123,123,123,213,231,312)$ & 720 & 62 & 47 & 15\\
43 & $(123,123,123,213,231,321)$ & 720 & 53 & 47 & 6\\
44 & $(123,123,123,213,312,312)$ & 360 & 53 & 47 & 6\\
45 & $(123,123,123,213,312,321)$ & 720 & 63 & 48 & 15\\
46 & $(123,123,123,213,321,321)$ & 360 & 54 & 48 & 6\\
47 & $(123,123,123,231,231,231)$ & 120 & 52 & 46 & 6\\
48 & $(123,123,123,231,231,312)$ & 360 & 64 & 47 & 17\\
49 & $(123,123,123,231,231,321)$ & 360 & 53 & 47 & 6\\
50 & $(123,123,123,231,312,312)$ & 360 & 64 & 47 & 17\\
51 & $(123,123,123,231,312,321)$ & 720 & 65 & 48 & 17\\
52 & $(123,123,123,231,321,321)$ & 360 & 54 & 48 & 6\\
53 & $(123,123,123,312,312,321)$ & 360 & 64 & 47 & 17\\
54 & $(123,123,123,312,321,321)$ & 360 & 65 & 48 & 17\\
55 & $(123,123,123,321,321,321)$ & 60 & 52 & 46 & 6\\
56 & $(123,123,132,132,213,213)$ & 540 & 52 & 46 & 6\\
57 & $(123,123,132,132,213,231)$ & 1080 & 64 & 47 & 17\\
58 & $(123,123,132,132,213,312)$ & 540 & 52 & 46 & 6\\
59 & $(123,123,132,132,213,321)$ & 1080 & 64 & 47 & 17\\
60 & $(123,123,132,132,231,231)$ & 540 & 67 & 48 & 19\\
61 & $(123,123,132,132,231,321)$ & 540 & 67 & 48 & 19\\
62 & $(123,123,132,213,213,231)$ & 540 & 61 & 46 & 15\\
63 & $(123,123,132,213,213,312)$ & 1080 & 53 & 47 & 6\\
64 & $(123,123,132,213,213,321)$ & 1080 & 62 & 47 & 15\\
65 & $(123,123,132,213,231,231)$ & 1080 & 64 & 47 & 17\\
66 & $(123,123,132,213,231,312)$ & 2160 & 65 & 48 & 17\\
67 & $(123,123,132,213,231,321)$ & 2160 & 65 & 48 & 17\\
68 & $(123,123,132,213,312,312)$ & 1080 & 53 & 47 & 6\\
69 & $(123,123,132,213,312,321)$ & 2160 & 65 & 48 & 17\\
70 & $(123,123,132,213,321,321)$ & 1080 & 65 & 48 & 17\\
71 & $(123,123,132,231,231,312)$ & 1080 & 67 & 48 & 19\\
72 & $(123,123,132,231,231,321)$ & 1080 & 65 & 48 & 17\\
73 & $(123,123,132,231,312,312)$ & 1080 & 64 & 47 & 17\\
74 & $(123,123,132,231,312,321)$ & 2160 & 67 & 48 & 19\\
75 & $(123,123,132,231,321,321)$ & 1080 & 65 & 48 & 17\\
76 & $(123,123,132,312,321,321)$ & 540 & 67 & 48 & 19\\
77 & $(123,123,213,213,312,312)$ & 540 & 54 & 48 & 6\\
78 & $(123,123,213,213,312,321)$ & 540 & 63 & 48 & 15\\
79 & $(123,123,213,231,312,312)$ & 1080 & 65 & 48 & 17\\
80 & $(123,123,213,231,312,321)$ & 2160 & 65 & 48 & 17\\
81 & $(123,123,213,231,321,321)$ & 540 & 54 & 48 & 6\\
82 & $(123,123,231,231,312,312)$ & 180 & 67 & 48 & 19\\
83 & $(123,132,213,231,312,321)$ & 720 & 67 & 48 & 19\\
\midrule
 & Total & 46656 & 4781 & 3865 & 916\\
\end{longtable}
}

\clearpage
{\small\setlength{\tabcolsep}{4.5pt}
\begin{longtable}{@{}r>{\raggedright\arraybackslash}p{12.9cm}@{}}
\caption{Dimension $ 6$: all additional facet inequalities beyond the baseline.}\label{tab:extras-n6}\\
\toprule
Case & Additional inequalities\\
\midrule\endfirsthead
\multicolumn{2}{l}{\small Table \ref{tab:extras-n6} (continued).}\\
\toprule
Case & Additional inequalities\\
\midrule\endhead
\midrule
\multicolumn{2}{r}{\small Continued on next page.}\\
\endfoot\bottomrule\endlastfoot
1 & None.\\[3pt]
2 & \(U_{23\mid 1}\leq \min\{t_{13},t_{23},t_{33},t_{43},t_{53}\}+d_{1}+d_{2}+d_{3}+d_{4}+d_{5}+d_{6}\)\\[3pt]
3 & None.\\[3pt]
4 & \(U_{13\mid 2}\leq \min\{t_{13},t_{23},t_{33},t_{43},t_{53}\}+d_{6}\)\\[3pt]
5 & \(U_{23\mid 1}\leq \min\{t_{13},t_{23},t_{33},t_{43},t_{53}\}+d_{1}+d_{2}+d_{3}+d_{4}+d_{5}\)\\[3pt]
6 & \(U_{13\mid 2}\leq \min\{t_{13},t_{23},t_{33},t_{43},t_{53}\}\)\\[3pt]
7 & \(U_{23\mid 1}\leq \min\{t_{13},t_{23},t_{33},t_{43},t_{53},t_{63}\}+d_{1}+d_{2}+d_{3}+d_{4}+d_{5}+d_{6}\)\\[3pt]
8 & \(U_{23\mid 1}\leq \min\{t_{13},t_{23},t_{33},t_{43},t_{63}\}+d_{1}+d_{2}+d_{3}+d_{4}+d_{5}\)\\[3pt]
9 & \(U_{23\mid 1}\leq \min\{t_{13},t_{23},t_{33},t_{43},t_{62}\}+d_{1}+d_{2}+d_{3}+d_{4}+d_{5}\)\newline \(U_{13\mid 2}\leq \min\{t_{13},t_{23},t_{33},t_{43},t_{52}\}+d_{6}\)\newline \(U_{12\mid 3}\leq \min\{t_{12},t_{22},t_{32},t_{42}\}\)\newline \(S\geq B_n\)\\[3pt]
10 & \(U_{23\mid 1}\leq \min\{t_{13},t_{23},t_{33},t_{43},t_{53},t_{63}\}+d_{1}+d_{2}+d_{3}+d_{4}+d_{5}\)\\[3pt]
11 & \(U_{23\mid 1}\leq \min\{t_{13},t_{23},t_{33},t_{43},t_{62}\}+d_{1}+d_{2}+d_{3}+d_{4}+d_{5}\)\newline \(U_{13\mid 2}\leq \min\{t_{13},t_{23},t_{33},t_{43},t_{52}\}\)\newline \(U_{12\mid 3}\leq \min\{t_{12},t_{22},t_{32},t_{42}\}+d_{6}\)\newline \(S\geq B_n\)\\[3pt]
12 & None.\\[3pt]
13 & \(U_{13\mid 2}\leq \min\{t_{13},t_{23},t_{33},t_{43},t_{53}\}+d_{5}+d_{6}\)\\[3pt]
14 & \(U_{23\mid 1}\leq \min\{t_{13},t_{23},t_{33},t_{43},t_{53}\}+d_{1}+d_{2}+d_{3}+d_{4}\)\\[3pt]
15 & \(U_{13\mid 2}\leq \min\{t_{13},t_{23},t_{33},t_{43},t_{53}\}+d_{5}\)\\[3pt]
16 & \(U_{13\mid 2}\leq \min\{t_{13},t_{23},t_{33},t_{43},t_{53},t_{63}\}+d_{5}+d_{6}\)\\[3pt]
17 & \(U_{23\mid 1}\leq \min\{t_{13},t_{23},t_{33},t_{43},t_{52}\}+d_{1}+d_{2}+d_{3}+d_{4}\)\newline \(U_{13\mid 2}\leq \min\{t_{13},t_{23},t_{33},t_{43},t_{62}\}+d_{5}\)\newline \(U_{12\mid 3}\leq \min\{t_{12},t_{22},t_{32},t_{42}\}+d_{6}\)\newline \(S\geq B_n\)\\[3pt]
18 & \(U_{13\mid 2}\leq \min\{t_{13},t_{23},t_{33},t_{43},t_{53},t_{63}\}+d_{5}\)\\[3pt]
19 & \(U_{23\mid 1}\leq \min\{t_{13},t_{23},t_{33},t_{43},t_{53},t_{63}\}+d_{1}+d_{2}+d_{3}+d_{4}\)\\[3pt]
20 & \(U_{23\mid 1}\leq \min\{t_{13},t_{23},t_{33},t_{43},t_{62}\}+d_{1}+d_{2}+d_{3}+d_{4}\)\newline \(U_{13\mid 2}\leq \min\{t_{13},t_{23},t_{33},t_{43},t_{52}\}\)\newline \(U_{12\mid 3}\leq \min\{t_{12},t_{22},t_{32},t_{42}\}+d_{5}+d_{6}\)\newline \(S\geq B_n\)\\[3pt]
21 & \(U_{13\mid 2}\leq \min\{t_{13},t_{23},t_{33},t_{43},t_{53},t_{63}\}\)\\[3pt]
22 & \(U_{23\mid 1}\leq \min\{t_{13},t_{23},t_{33},t_{43},t_{53},t_{63}\}+d_{1}+d_{2}+d_{3}+d_{4}+d_{5}+d_{6}\)\\[3pt]
23 & \(U_{23\mid 1}\leq \min\{t_{13},t_{23},t_{33},t_{43},t_{53},t_{63}\}+d_{1}+d_{2}+d_{3}+d_{4}+d_{5}\)\\[3pt]
24 & \(U_{23\mid 1}\leq \min\{t_{13},t_{23},t_{33},t_{43},t_{53},t_{62}\}+d_{1}+d_{2}+d_{3}+d_{4}+d_{5}\)\newline \(U_{13\mid 2}\leq \min\{t_{13},t_{23},t_{33},t_{42},t_{52}\}+d_{6}\)\newline \(U_{12\mid 3}\leq \min\{t_{12},t_{22},t_{32},t_{43},t_{53}\}\)\newline \(S\geq B_n\)\\[3pt]
25 & \(U_{23\mid 1}\leq \min\{t_{13},t_{23},t_{33},t_{43},t_{53},t_{63}\}+d_{1}+d_{2}+d_{3}+d_{4}+d_{5}\)\\[3pt]
26 & \(U_{23\mid 1}\leq \min\{t_{13},t_{23},t_{33},t_{43},t_{53},t_{62}\}+d_{1}+d_{2}+d_{3}+d_{4}+d_{5}\)\newline \(U_{13\mid 2}\leq \min\{t_{13},t_{23},t_{33},t_{42},t_{52}\}\)\newline \(U_{12\mid 3}\leq \min\{t_{12},t_{22},t_{32},t_{43},t_{53}\}+d_{6}\)\newline \(S\geq B_n\)\\[3pt]
27 & \(U_{23\mid 1}\leq \min\{t_{13},t_{23},t_{33},t_{53},t_{63}\}+d_{1}+d_{2}+d_{3}+d_{4}\)\\[3pt]
28 & \(U_{23\mid 1}\leq \min\{t_{13},t_{23},t_{33},t_{53},t_{62}\}+d_{1}+d_{2}+d_{3}+d_{4}\)\newline \(U_{13\mid 2}\leq \min\{t_{13},t_{23},t_{33},t_{42},t_{53}\}+d_{5}+d_{6}\)\newline \(U_{12\mid 3}\leq \min\{t_{12},t_{22},t_{32},t_{52}\}\)\newline \(S\geq B_n\)\\[3pt]
29 & \(U_{23\mid 1}\leq \min\{t_{13},t_{23},t_{33},t_{43},t_{53},t_{63}\}+d_{1}+d_{2}+d_{3}+d_{4}\)\\[3pt]
30 & \(U_{23\mid 1}\leq \min\{t_{13},t_{23},t_{33},t_{53},t_{62}\}+d_{1}+d_{2}+d_{3}+d_{4}\)\newline \(U_{13\mid 2}\leq \min\{t_{13},t_{23},t_{33},t_{42},t_{53}\}+d_{5}\)\newline \(U_{12\mid 3}\leq \min\{t_{12},t_{22},t_{32},t_{52}\}+d_{6}\)\newline \(S\geq B_n\)\\[3pt]
31 & \(U_{23\mid 1}\leq \min\{t_{13},t_{23},t_{33},t_{52},t_{62}\}+d_{1}+d_{2}+d_{3}+d_{4}\)\newline \(U_{13\mid 2}\leq \min\{t_{13},t_{23},t_{33},t_{42},t_{53},t_{63}\}+d_{5}+d_{6}\)\newline \(U_{12\mid 3}\leq \min\{t_{12},t_{22},t_{32},t_{53},t_{63}\}\)\newline \(S\geq B_n\)\\[3pt]
32 & \(U_{23\mid 1}\leq \min\{t_{13},t_{23},t_{33},t_{43},t_{52},t_{63}\}+d_{1}+d_{2}+d_{3}+d_{4}\)\newline \(U_{13\mid 2}\leq \min\{t_{13},t_{23},t_{33},t_{42},t_{62}\}+d_{5}\)\newline \(U_{12\mid 3}\leq \min\{t_{12},t_{22},t_{32},t_{43},t_{63}\}+d_{6}\)\newline \(S\geq B_n\)\\[3pt]
33 & \(U_{23\mid 1}\leq \min\{t_{13},t_{23},t_{33},t_{52},t_{62}\}+d_{1}+d_{2}+d_{3}+d_{4}\)\newline \(U_{13\mid 2}\leq \min\{t_{13},t_{23},t_{33},t_{42},t_{53},t_{63}\}+d_{5}\)\newline \(U_{12\mid 3}\leq \min\{t_{12},t_{22},t_{32},t_{53},t_{63}\}+d_{6}\)\newline \(S\geq B_n\)\\[3pt]
34 & \(U_{23\mid 1}\leq \min\{t_{13},t_{23},t_{33},t_{43},t_{53},t_{63}\}+d_{1}+d_{2}+d_{3}+d_{4}\)\\[3pt]
35 & \(U_{23\mid 1}\leq \min\{t_{13},t_{23},t_{33},t_{43},t_{53},t_{62}\}+d_{1}+d_{2}+d_{3}+d_{4}\)\newline \(U_{13\mid 2}\leq \min\{t_{13},t_{23},t_{33},t_{42},t_{52}\}\)\newline \(U_{12\mid 3}\leq \min\{t_{12},t_{22},t_{32},t_{43},t_{53}\}+d_{5}+d_{6}\)\newline \(S\geq B_n\)\\[3pt]
36 & \(U_{23\mid 1}\leq \min\{t_{13},t_{23},t_{33},t_{52},t_{62}\}+d_{1}+d_{2}+d_{3}+d_{4}\)\newline \(U_{13\mid 2}\leq \min\{t_{13},t_{23},t_{33},t_{42},t_{53},t_{63}\}\)\newline \(U_{12\mid 3}\leq \min\{t_{12},t_{22},t_{32},t_{53},t_{63}\}+d_{5}+d_{6}\)\newline \(S\geq B_n\)\\[3pt]
37 & None.\\[3pt]
38 & \(U_{13\mid 2}\leq \min\{t_{13},t_{23},t_{33},t_{43},t_{53}\}+d_{4}+d_{5}+d_{6}\)\\[3pt]
39 & \(U_{23\mid 1}\leq \min\{t_{13},t_{23},t_{33},t_{43},t_{53}\}+d_{1}+d_{2}+d_{3}\)\\[3pt]
40 & \(U_{13\mid 2}\leq \min\{t_{13},t_{23},t_{33},t_{43},t_{53}\}+d_{4}+d_{5}\)\\[3pt]
41 & \(U_{13\mid 2}\leq \min\{t_{13},t_{23},t_{33},t_{43},t_{53},t_{63}\}+d_{4}+d_{5}+d_{6}\)\\[3pt]
42 & \(U_{23\mid 1}\leq \min\{t_{13},t_{23},t_{33},t_{43},t_{52}\}+d_{1}+d_{2}+d_{3}\)\newline \(U_{13\mid 2}\leq \min\{t_{13},t_{23},t_{33},t_{43},t_{62}\}+d_{4}+d_{5}\)\newline \(U_{12\mid 3}\leq \min\{t_{12},t_{22},t_{32},t_{42}\}+d_{6}\)\newline \(S\geq B_n\)\\[3pt]
43 & \(U_{13\mid 2}\leq \min\{t_{13},t_{23},t_{33},t_{43},t_{53},t_{63}\}+d_{4}+d_{5}\)\\[3pt]
44 & \(U_{23\mid 1}\leq \min\{t_{13},t_{23},t_{33},t_{43},t_{53},t_{63}\}+d_{1}+d_{2}+d_{3}\)\\[3pt]
45 & \(U_{23\mid 1}\leq \min\{t_{13},t_{23},t_{33},t_{43},t_{62}\}+d_{1}+d_{2}+d_{3}\)\newline \(U_{13\mid 2}\leq \min\{t_{13},t_{23},t_{33},t_{43},t_{52}\}+d_{4}\)\newline \(U_{12\mid 3}\leq \min\{t_{12},t_{22},t_{32},t_{42}\}+d_{5}+d_{6}\)\newline \(S\geq B_n\)\\[3pt]
46 & \(U_{13\mid 2}\leq \min\{t_{13},t_{23},t_{33},t_{43},t_{53},t_{63}\}+d_{4}\)\\[3pt]
47 & \(U_{13\mid 2}\leq \min\{t_{13},t_{23},t_{33},t_{43},t_{53},t_{63}\}+d_{4}+d_{5}+d_{6}\)\\[3pt]
48 & \(U_{23\mid 1}\leq \min\{t_{13},t_{23},t_{33},t_{42},t_{52}\}+d_{1}+d_{2}+d_{3}\)\newline \(U_{13\mid 2}\leq \min\{t_{13},t_{23},t_{33},t_{43},t_{53},t_{62}\}+d_{4}+d_{5}\)\newline \(U_{12\mid 3}\leq \min\{t_{12},t_{22},t_{32},t_{43},t_{53}\}+d_{6}\)\newline \(S\geq B_n\)\\[3pt]
49 & \(U_{13\mid 2}\leq \min\{t_{13},t_{23},t_{33},t_{43},t_{53},t_{63}\}+d_{4}+d_{5}\)\\[3pt]
50 & \(U_{23\mid 1}\leq \min\{t_{13},t_{23},t_{33},t_{42},t_{53},t_{63}\}+d_{1}+d_{2}+d_{3}\)\newline \(U_{13\mid 2}\leq \min\{t_{13},t_{23},t_{33},t_{52},t_{62}\}+d_{4}\)\newline \(U_{12\mid 3}\leq \min\{t_{12},t_{22},t_{32},t_{53},t_{63}\}+d_{5}+d_{6}\)\newline \(S\geq B_n\)\\[3pt]
51 & \(U_{23\mid 1}\leq \min\{t_{13},t_{23},t_{33},t_{42},t_{62}\}+d_{1}+d_{2}+d_{3}\)\newline \(U_{13\mid 2}\leq \min\{t_{13},t_{23},t_{33},t_{43},t_{52},t_{63}\}+d_{4}\)\newline \(U_{12\mid 3}\leq \min\{t_{12},t_{22},t_{32},t_{43},t_{63}\}+d_{5}+d_{6}\)\newline \(S\geq B_n\)\\[3pt]
52 & \(U_{13\mid 2}\leq \min\{t_{13},t_{23},t_{33},t_{43},t_{53},t_{63}\}+d_{4}\)\\[3pt]
53 & \(U_{23\mid 1}\leq \min\{t_{13},t_{23},t_{33},t_{43},t_{53},t_{62}\}+d_{1}+d_{2}+d_{3}\)\newline \(U_{13\mid 2}\leq \min\{t_{13},t_{23},t_{33},t_{42},t_{52}\}\)\newline \(U_{12\mid 3}\leq \min\{t_{12},t_{22},t_{32},t_{43},t_{53}\}+d_{4}+d_{5}+d_{6}\)\newline \(S\geq B_n\)\\[3pt]
54 & \(U_{23\mid 1}\leq \min\{t_{13},t_{23},t_{33},t_{52},t_{62}\}+d_{1}+d_{2}+d_{3}\)\newline \(U_{13\mid 2}\leq \min\{t_{13},t_{23},t_{33},t_{42},t_{53},t_{63}\}\)\newline \(U_{12\mid 3}\leq \min\{t_{12},t_{22},t_{32},t_{53},t_{63}\}+d_{4}+d_{5}+d_{6}\)\newline \(S\geq B_n\)\\[3pt]
55 & \(U_{13\mid 2}\leq \min\{t_{13},t_{23},t_{33},t_{43},t_{53},t_{63}\}\)\\[3pt]
56 & \(U_{23\mid 1}\leq \min\{t_{13},t_{23},t_{33},t_{43},t_{53},t_{63}\}+d_{1}+d_{2}+d_{3}+d_{4}\)\\[3pt]
57 & \(U_{23\mid 1}\leq \min\{t_{13},t_{23},t_{33},t_{43},t_{53},t_{62}\}+d_{1}+d_{2}+d_{3}+d_{4}\)\newline \(U_{13\mid 2}\leq \min\{t_{13},t_{23},t_{32},t_{42},t_{53}\}+d_{5}+d_{6}\)\newline \(U_{12\mid 3}\leq \min\{t_{12},t_{22},t_{33},t_{43},t_{52}\}\)\newline \(S\geq B_n\)\\[3pt]
58 & \(U_{23\mid 1}\leq \min\{t_{13},t_{23},t_{33},t_{43},t_{53},t_{63}\}+d_{1}+d_{2}+d_{3}+d_{4}\)\\[3pt]
59 & \(U_{23\mid 1}\leq \min\{t_{13},t_{23},t_{33},t_{43},t_{53},t_{62}\}+d_{1}+d_{2}+d_{3}+d_{4}\)\newline \(U_{13\mid 2}\leq \min\{t_{13},t_{23},t_{32},t_{42},t_{53}\}+d_{5}\)\newline \(U_{12\mid 3}\leq \min\{t_{12},t_{22},t_{33},t_{43},t_{52}\}+d_{6}\)\newline \(S\geq B_n\)\\[3pt]
60 & \(U_{23\mid 1}\leq \min\{t_{13},t_{23},t_{33},t_{43},t_{52},t_{62}\}+d_{1}+d_{2}+d_{3}+d_{4}\)\newline \(U_{13\mid 2}\leq \min\{t_{13},t_{23},t_{32},t_{42},t_{53},t_{63}\}+d_{5}+d_{6}\)\newline \(U_{12\mid 3}\leq \min\{t_{12},t_{22},t_{33},t_{43},t_{53},t_{63}\}\)\newline \(S\geq B_n\)\\[3pt]
61 & \(U_{23\mid 1}\leq \min\{t_{13},t_{23},t_{33},t_{43},t_{52},t_{62}\}+d_{1}+d_{2}+d_{3}+d_{4}\)\newline \(U_{13\mid 2}\leq \min\{t_{13},t_{23},t_{32},t_{42},t_{53},t_{63}\}+d_{5}\)\newline \(U_{12\mid 3}\leq \min\{t_{12},t_{22},t_{33},t_{43},t_{53},t_{63}\}+d_{6}\)\newline \(S\geq B_n\)\\[3pt]
62 & \(U_{23\mid 1}\leq \min\{t_{13},t_{23},t_{43},t_{53},t_{62}\}+d_{1}+d_{2}+d_{3}\)\newline \(U_{13\mid 2}\leq \min\{t_{13},t_{23},t_{32},t_{43},t_{53}\}+d_{4}+d_{5}+d_{6}\)\newline \(U_{12\mid 3}\leq \min\{t_{12},t_{22},t_{42},t_{52}\}\)\newline \(S\geq B_n\)\\[3pt]
63 & \(U_{23\mid 1}\leq \min\{t_{13},t_{23},t_{33},t_{43},t_{53},t_{63}\}+d_{1}+d_{2}+d_{3}\)\\[3pt]
64 & \(U_{23\mid 1}\leq \min\{t_{13},t_{23},t_{43},t_{53},t_{62}\}+d_{1}+d_{2}+d_{3}\)\newline \(U_{13\mid 2}\leq \min\{t_{13},t_{23},t_{32},t_{43},t_{53}\}+d_{4}+d_{5}\)\newline \(U_{12\mid 3}\leq \min\{t_{12},t_{22},t_{42},t_{52}\}+d_{6}\)\newline \(S\geq B_n\)\\[3pt]
65 & \(U_{23\mid 1}\leq \min\{t_{13},t_{23},t_{43},t_{52},t_{62}\}+d_{1}+d_{2}+d_{3}\)\newline \(U_{13\mid 2}\leq \min\{t_{13},t_{23},t_{32},t_{43},t_{53},t_{63}\}+d_{4}+d_{5}+d_{6}\)\newline \(U_{12\mid 3}\leq \min\{t_{12},t_{22},t_{42},t_{53},t_{63}\}\)\newline \(S\geq B_n\)\\[3pt]
66 & \(U_{23\mid 1}\leq \min\{t_{13},t_{23},t_{33},t_{43},t_{52},t_{63}\}+d_{1}+d_{2}+d_{3}\)\newline \(U_{13\mid 2}\leq \min\{t_{13},t_{23},t_{32},t_{43},t_{62}\}+d_{4}+d_{5}\)\newline \(U_{12\mid 3}\leq \min\{t_{12},t_{22},t_{33},t_{42},t_{63}\}+d_{6}\)\newline \(S\geq B_n\)\\[3pt]
67 & \(U_{23\mid 1}\leq \min\{t_{13},t_{23},t_{43},t_{52},t_{62}\}+d_{1}+d_{2}+d_{3}\)\newline \(U_{13\mid 2}\leq \min\{t_{13},t_{23},t_{32},t_{43},t_{53},t_{63}\}+d_{4}+d_{5}\)\newline \(U_{12\mid 3}\leq \min\{t_{12},t_{22},t_{42},t_{53},t_{63}\}+d_{6}\)\newline \(S\geq B_n\)\\[3pt]
68 & \(U_{23\mid 1}\leq \min\{t_{13},t_{23},t_{33},t_{43},t_{53},t_{63}\}+d_{1}+d_{2}+d_{3}\)\\[3pt]
69 & \(U_{23\mid 1}\leq \min\{t_{13},t_{23},t_{33},t_{43},t_{53},t_{62}\}+d_{1}+d_{2}+d_{3}\)\newline \(U_{13\mid 2}\leq \min\{t_{13},t_{23},t_{32},t_{43},t_{52}\}+d_{4}\)\newline \(U_{12\mid 3}\leq \min\{t_{12},t_{22},t_{33},t_{42},t_{53}\}+d_{5}+d_{6}\)\newline \(S\geq B_n\)\\[3pt]
70 & \(U_{23\mid 1}\leq \min\{t_{13},t_{23},t_{43},t_{52},t_{62}\}+d_{1}+d_{2}+d_{3}\)\newline \(U_{13\mid 2}\leq \min\{t_{13},t_{23},t_{32},t_{43},t_{53},t_{63}\}+d_{4}\)\newline \(U_{12\mid 3}\leq \min\{t_{12},t_{22},t_{42},t_{53},t_{63}\}+d_{5}+d_{6}\)\newline \(S\geq B_n\)\\[3pt]
71 & \(U_{23\mid 1}\leq \min\{t_{13},t_{23},t_{33},t_{42},t_{52},t_{63}\}+d_{1}+d_{2}+d_{3}\)\newline \(U_{13\mid 2}\leq \min\{t_{13},t_{23},t_{32},t_{43},t_{53},t_{62}\}+d_{4}+d_{5}\)\newline \(U_{12\mid 3}\leq \min\{t_{12},t_{22},t_{33},t_{43},t_{53},t_{63}\}+d_{6}\)\newline \(S\geq B_n\)\\[3pt]
72 & \(U_{23\mid 1}\leq \min\{t_{13},t_{23},t_{42},t_{52},t_{62}\}+d_{1}+d_{2}+d_{3}\)\newline \(U_{13\mid 2}\leq \min\{t_{13},t_{23},t_{32},t_{43},t_{53},t_{63}\}+d_{4}+d_{5}\)\newline \(U_{12\mid 3}\leq \min\{t_{12},t_{22},t_{43},t_{53},t_{63}\}+d_{6}\)\newline \(S\geq B_n\)\\[3pt]
73 & \(U_{23\mid 1}\leq \min\{t_{13},t_{23},t_{33},t_{42},t_{53},t_{63}\}+d_{1}+d_{2}+d_{3}\)\newline \(U_{13\mid 2}\leq \min\{t_{13},t_{23},t_{32},t_{52},t_{62}\}+d_{4}\)\newline \(U_{12\mid 3}\leq \min\{t_{12},t_{22},t_{33},t_{53},t_{63}\}+d_{5}+d_{6}\)\newline \(S\geq B_n\)\\[3pt]
74 & \(U_{23\mid 1}\leq \min\{t_{13},t_{23},t_{33},t_{42},t_{53},t_{62}\}+d_{1}+d_{2}+d_{3}\)\newline \(U_{13\mid 2}\leq \min\{t_{13},t_{23},t_{32},t_{43},t_{52},t_{63}\}+d_{4}\)\newline \(U_{12\mid 3}\leq \min\{t_{12},t_{22},t_{33},t_{43},t_{53},t_{63}\}+d_{5}+d_{6}\)\newline \(S\geq B_n\)\\[3pt]
75 & \(U_{23\mid 1}\leq \min\{t_{13},t_{23},t_{42},t_{52},t_{62}\}+d_{1}+d_{2}+d_{3}\)\newline \(U_{13\mid 2}\leq \min\{t_{13},t_{23},t_{32},t_{43},t_{53},t_{63}\}+d_{4}\)\newline \(U_{12\mid 3}\leq \min\{t_{12},t_{22},t_{43},t_{53},t_{63}\}+d_{5}+d_{6}\)\newline \(S\geq B_n\)\\[3pt]
76 & \(U_{23\mid 1}\leq \min\{t_{13},t_{23},t_{33},t_{43},t_{52},t_{62}\}+d_{1}+d_{2}+d_{3}\)\newline \(U_{13\mid 2}\leq \min\{t_{13},t_{23},t_{32},t_{42},t_{53},t_{63}\}\)\newline \(U_{12\mid 3}\leq \min\{t_{12},t_{22},t_{33},t_{43},t_{53},t_{63}\}+d_{4}+d_{5}+d_{6}\)\newline \(S\geq B_n\)\\[3pt]
77 & \(U_{23\mid 1}\leq \min\{t_{13},t_{23},t_{33},t_{43},t_{53},t_{63}\}+d_{1}+d_{2}\)\\[3pt]
78 & \(U_{23\mid 1}\leq \min\{t_{13},t_{23},t_{33},t_{43},t_{62}\}+d_{1}+d_{2}\)\newline \(U_{13\mid 2}\leq \min\{t_{13},t_{23},t_{33},t_{43},t_{52}\}+d_{3}+d_{4}\)\newline \(U_{12\mid 3}\leq \min\{t_{12},t_{22},t_{32},t_{42}\}+d_{5}+d_{6}\)\newline \(S\geq B_n\)\\[3pt]
79 & \(U_{23\mid 1}\leq \min\{t_{13},t_{23},t_{33},t_{42},t_{53},t_{63}\}+d_{1}+d_{2}\)\newline \(U_{13\mid 2}\leq \min\{t_{13},t_{23},t_{33},t_{52},t_{62}\}+d_{3}+d_{4}\)\newline \(U_{12\mid 3}\leq \min\{t_{12},t_{22},t_{32},t_{53},t_{63}\}+d_{5}+d_{6}\)\newline \(S\geq B_n\)\\[3pt]
80 & \(U_{23\mid 1}\leq \min\{t_{13},t_{23},t_{33},t_{42},t_{62}\}+d_{1}+d_{2}\)\newline \(U_{13\mid 2}\leq \min\{t_{13},t_{23},t_{33},t_{43},t_{52},t_{63}\}+d_{3}+d_{4}\)\newline \(U_{12\mid 3}\leq \min\{t_{12},t_{22},t_{32},t_{43},t_{63}\}+d_{5}+d_{6}\)\newline \(S\geq B_n\)\\[3pt]
81 & \(U_{13\mid 2}\leq \min\{t_{13},t_{23},t_{33},t_{43},t_{53},t_{63}\}+d_{3}+d_{4}\)\\[3pt]
82 & \(U_{23\mid 1}\leq \min\{t_{13},t_{23},t_{32},t_{42},t_{53},t_{63}\}+d_{1}+d_{2}\)\newline \(U_{13\mid 2}\leq \min\{t_{13},t_{23},t_{33},t_{43},t_{52},t_{62}\}+d_{3}+d_{4}\)\newline \(U_{12\mid 3}\leq \min\{t_{12},t_{22},t_{33},t_{43},t_{53},t_{63}\}+d_{5}+d_{6}\)\newline \(S\geq B_n\)\\[3pt]
83 & \(U_{23\mid 1}\leq \min\{t_{13},t_{23},t_{33},t_{42},t_{53},t_{62}\}+d_{1}+d_{2}\)\newline \(U_{13\mid 2}\leq \min\{t_{13},t_{22},t_{33},t_{43},t_{52},t_{63}\}+d_{3}+d_{4}\)\newline \(U_{12\mid 3}\leq \min\{t_{12},t_{23},t_{32},t_{43},t_{53},t_{63}\}+d_{5}+d_{6}\)\newline \(S\geq B_n\)\\[3pt]
\end{longtable}
}

% END GENERATED APPENDIX TABLES
\clearpage

\section{Two explicit two-point checkerboards in dimension three}\label{app:two-examples}
We record a comparable and an incomparable example of the finite mass
construction in \cref{lem:grid}. In both, $M_k=(m_{abk})_{a,b=1}^3$: rows index the first
coordinate, columns the second, and the slice index $k$ the third.
All displayed matrices contain \emph{cell masses}. Dividing an entry by
the product of its three band widths gives the density on that cell.

\subsection{Comparable points}
Let
\[
 c=(0.4,0.5,0.6),\quad d=(0.7,0.8,0.9),\qquad u=0.2,\quad v=0.5.
\]
The arrays below define copulas $C_-,C_+$ with
\[
C_-(c)=C_+(c)=0.2,\qquad C_-(d)=0.4,\qquad C_+(d)=0.7.
\]
The desired value $v=0.5$ is obtained with mixture weight
$(0.5-0.4)/(0.7-0.4)=1/3$.
The three bands in each coordinate are cut at $(c_i,d_i)$, and their
width vectors are $(0.4,0.3,0.3)$, $(0.5,0.3,0.2)$, and $(0.6,0.3,0.1)$.
The two copulas are specified by the following cell masses:
\[\begin{aligned}
M^{-}_{1}&=\frac1{10}\mat{2&0&1\\0&0&1\\2&0&0}&
M^{-}_{2}&=\frac1{10}\mat{0&0&0\\0&2&0\\0&1&0}&
M^{-}_{3}&=\frac1{10}\mat{1&0&0\\0&0&0\\0&0&0}\\[5pt]
M^{+}_{1}&=\frac1{10}\mat{2&1&0\\0&1&0\\0&0&2}&
M^{+}_{2}&=\frac1{10}\mat{0&1&0\\2&0&0\\0&0&0}&
M^{+}_{3}&=\frac1{10}\mat{0&0&0\\0&0&0\\1&0&0}.
\end{aligned}\]
For example, the lower endpoint comes from the following membership
partition of the auxiliary interval:
\begin{center}\small\setlength{\tabcolsep}{3.6pt}
\begin{tabular}{c@{\quad}ccccccc}
\toprule
$\omega$ & $[0,.1)$ & $[.1,.3)$ & $[.3,.4)$ & $[.4,.5)$ & $[.5,.6)$ & $[.6,.8)$ & $[.8,1)$\\
Cell & $322$ & $311$ & $231$ & $131$ & $113$ & $222$ & $111$\\
\bottomrule
\end{tabular}
\end{center}
Define $D_{ij}$ as the union of intervals whose cell has $i$th index $j$;
then $A_i=D_{i1}$ and $B_i=D_{i1}\cup D_{i2}$. This gives
$\bigcap_iA_i=[.8,1)$ and $\bigcap_iB_i=[.6,1)$, of measures $0.2$ and
$0.4$. For instance $[.6,.8)$ contributes mass $0.2$ to cell $222$,
whose volume is $0.3^3$ and whose density is $200/27$.

The mixture $\tfrac23C_-+\tfrac13C_+$ has masses
\[\begin{aligned}
M_1&=\frac1{30}\mat{6&1&2\\0&1&2\\4&0&2}&
M_2&=\frac1{30}\mat{0&1&0\\2&4&0\\0&2&0}&
M_3&=\frac1{30}\mat{2&0&0\\0&0&0\\1&0&0}.
\end{aligned}\]
Each endpoint and the mixture have the stated marginal width vectors.
For the mixture, $C(c)=m_{111}=6/30=0.2$ and
$C(d)=\sum_{a,b,k\leq2}m_{abk}=15/30=0.5$.

\subsection{Incomparable points}
Let
\[
 c=(0.7,0.6,0.5),\quad d=(0.5,0.8,0.7),\qquad u=0.4,\quad v=0.35.
\]
The two arrays below define copulas with
\[
C_-(c)=C_+(c)=0.4,\qquad C_-(d)=0.2,\qquad C_+(d)=0.5.
\]
Their average has the desired value $v=0.35$ at $d$.
The cuts are the smaller and larger of $c_i,d_i$; the width vectors are
$(0.5,0.2,0.3)$, $(0.6,0.2,0.2)$, and $(0.5,0.2,0.3)$.
Endpoint mass arrays are
\[\begin{aligned}
M^{-}_{1}&=\frac1{10}\mat{2&0&1\\2&0&0\\0&0&0}&
M^{-}_{2}&=\frac1{10}\mat{0&0&1\\0&0&0\\1&0&0}&
M^{-}_{3}&=\frac1{10}\mat{1&0&0\\0&0&0\\0&2&0}\\[5pt]
M^{+}_{1}&=\frac1{10}\mat{4&0&0\\0&1&0\\0&0&0}&
M^{+}_{2}&=\frac1{10}\mat{1&0&0\\0&1&0\\0&0&0}&
M^{+}_{3}&=\frac1{10}\mat{0&0&0\\0&0&0\\1&0&2}.
\end{aligned}\]
The lower endpoint can be read from the partition
\begin{center}\small\setlength{\tabcolsep}{3.6pt}
\begin{tabular}{c@{\quad}ccccccc}
\toprule
$\omega$ & $[0,.2)$ & $[.2,.3)$ & $[.3,.4)$ & $[.4,.5)$ & $[.5,.6)$ & $[.6,.8)$ & $[.8,1)$\\
Cell & $323$ & $312$ & $132$ & $131$ & $113$ & $211$ & $111$\\
\bottomrule
\end{tabular}
\end{center}
Here $A_1=D_{11}\cup D_{12}$ and $B_1=D_{11}$, whereas
$A_i=D_{i1}$ and $B_i=D_{i1}\cup D_{i2}$ for $i=2,3$.
The interval $[.6,.8)$ lies in all $A_i$ but outside $B_1$, accounting
for the difference $C_-(c)-C_-(d)=0.2$.
The average of the endpoint arrays is
\[\begin{aligned}
M_1&=\frac1{20}\mat{6&0&1\\2&1&0\\0&0&0}&
M_2&=\frac1{20}\mat{1&0&1\\0&1&0\\1&0&0}&
M_3&=\frac1{20}\mat{1&0&0\\0&0&0\\1&2&2}.
\end{aligned}\]
The marginal sums agree with the band widths, and
\[
 C(c)=m_{111}+m_{211}=8/20=0.4,\qquad
 C(d)=\sum_{b,k\leq2}m_{1bk}=7/20=0.35.
\]
Together with \eqref{eq:checker-formula}, these arrays specify the full
interpolating copula, not only its values at $c$ and $d$.


\section{Reproducibility of the exact calculations}\label{app:verification}
The supplement contains four certificates and their verification programs,
together with rational checks of the examples. With Python 3.10 or later,
run the following commands from the supplement directory:
\begin{verbatim}
python3 verify_three_point.py
python3 verify_all.py
python3 verify_appendix_tables.py --recompute 4 5
python3 verify_examples.py
python3 check_expanded_examples.py
\end{verbatim}
Only the Python standard library is required. Run without \texttt{-O},
because verification uses assertions. The second command runs
\texttt{verify\_base.py} and \texttt{verify\_split.py}; these can also be
run separately. Reconstructing the dimension-six base can take several
minutes. The verification uses exact integer and rational arithmetic.

\paragraph{The trivariate facet tables.}
The program \texttt{verify\_three\_point.py} generates all $216$ order
triples, reduces them to the ten representatives, constructs the $64$
vertices for each, and enumerates every facet from scratch. It reconstructs
the inequality families and both trivariate tables in
\cref{sec:configurations}. The certificate
\texttt{facet\_certificate.json} uses the coordinate order \eqref{eq:z},
followed by the constant entry $1$. Each stored normal $a$ represents
$a\cdot(\xi,1)\geq0$. Across the ten representatives it records $283$
primitive facet normals, as well as the rank matrices, orbit sizes, and
facet counts. Completeness is checked through the successive
half-space-intersection algorithm, rather than validity alone.

\paragraph{The finite bases and the projection step.}
The program \texttt{verify\_base.py} reconstructs the full facet lists in
dimensions two and six and compares them with
\texttt{base\_certificate.json}. The program
\texttt{verify\_split.py} reconstructs all positive circuits and all
projected inequalities used in \cref{lem:merge-selected}; it then checks
all $11{,}358$ nonnegative rational identities stored in
\texttt{split\_certificate.json}. The projection calculation permits
arbitrary nonnegative totals for the remaining coordinate bands. Thus
it establishes a dimension-independent implication, not a numerical test
in selected higher dimensions. Dimensions above six follow from the
merging lemma and induction; they are not separately enumerated.

\paragraph{The complete appendix tables.}
The certificate \texttt{appendix\_certificate.json} contains all facet
normals for dimensions two through six. The program
\texttt{verify\_appendix\_tables.py} checks every representative and orbit
size, reconstructs the facet counts, and expands every printed minimum
into its individual linear inequalities. Their primitive normals must
coincide exactly with the facets outside the baseline. It also regenerates
\texttt{appendix\_tables.tex}, whose contents are included in the standalone
manuscript source. With \texttt{--recompute 4 5}, the program independently
reconstructs all facets in dimensions four and five; the preceding
commands reconstruct dimensions two, three, and six. Thus every dimension
listed in Appendix~\ref{app:all-tables} is covered by a complete exact
enumeration. One may also use \texttt{--recompute 2 3 4 5 6} to reconstruct
all five dimensions in this program alone.

\paragraph{The examples.}
The program \texttt{verify\_examples.py} checks the rational checkerboard
examples in \cref{sec:examples,app:two-examples}.
The program \texttt{check\_expanded\_examples.py} checks the two worked
eliminations against the exact row systems for all six orders of the
selected column and verifies the numerical merging illustration using
fractions. Verification logs and a README are supplied with the source.

\begin{thebibliography}{9}
\bibitem{DeBaets2013}
B.~De Baets, H.~De Meyer, J.~Fern\'andez-S\'anchez, and M.~\'Ubeda-Flores,
\emph{On the existence of a trivariate copula with given values of a
trivariate quasi-copula at several points},
Fuzzy Sets and Systems \textbf{228} (2013), 3--14.
\href{https://doi.org/10.1016/j.fss.2012.07.006}{doi:10.1016/j.fss.2012.07.006}.

\bibitem{FukudaProdon1996}
K.~Fukuda and A.~Prodon,
\emph{Double description method revisited},
in Combinatorics and Computer Science,
Lecture Notes in Computer Science, vol.~1120,
Springer, Berlin, 1996, pp.~91--111.

\bibitem{Klement2023}
E.~P.~Klement, D.~Kokol Bukov\v{s}ek, M.~Omladi\v{c},
S.~Saminger-Platz, and N.~Stopar,
\emph{Multivariate copulas with given values at two arbitrary points},
Statistical Papers \textbf{64} (2023), 2015--2055.
\href{https://doi.org/10.1007/s00362-022-01362-4}{doi:10.1007/s00362-022-01362-4}.

\bibitem{Schrijver1986}
A.~Schrijver,
\emph{Theory of Linear and Integer Programming},
John Wiley \& Sons, Chichester, 1986.

\bibitem{OVZ}
M.~Omladi\v{c}, M.~Vuk, and A.~Zalar,
\emph{A linear programming approach to finite copula interpolation},
\url{https://zalara.github.io/Papers/linear_programming_approach_to_finite_copula_interpolation.pdf}.
\end{thebibliography}
\end{document}
